%TCIDATA{LaTeXparent=0,0,EOMN.tex}
                      

\section{Dual Operator Spaces}

\subsection{Operators mapping ${\cal B}_{1}\left( {\cal H}^{N}\right)
\longrightarrow {\cal B}_{1}\left( {\cal H}^{N}\right) $}

The dual space of the Banach space, ${\cal B}_{1}\left( {\cal H}^{N}\right) $%
, of operators formed by the set of trace class operators equipped with the
norm%
\begin{equation}
\left\| Y\right\| _{1}=Tr\left\{ \left( Y^{\dagger }Y\right) ^{\frac{1}{2}%
}\right\}
\end{equation}%
is the Banach space of bounded operators ${\cal B}\left( {\cal H}^{N}\right) 
$ equipped with the norm%
\begin{equation}
\left\| X\right\| =%
%TCIMACRO{
%\underset{\left\| \left| \Psi \right\rangle \right\| _{{\cal H}^{N}}=1}{\sup }}%
%BeginExpansion
\mathrel{\mathop{\sup }\limits_{\left\| \left| \Psi \right\rangle \right\| _{{\cal H}^{N}}=1}}%
%EndExpansion
\left\| X\left| \Psi \right\rangle \right\| _{{\cal H}^{N}};\quad \left|
\Psi \right\rangle \in {\cal H}^{N}
\end{equation}%
where $\left\| \left| \Psi \right\rangle \right\| _{{\cal H}^{N}}=\left(
\left\langle \Psi |\Psi \right\rangle \right) ^{\frac{1}{2}}$. The duality
is expressed by%
\begin{equation}
f_{X}\left( Y\right) =Tr\left\{ \left( X^{\dagger }Y\right) \right\} \equiv
\left( X|Y\right) \quad \forall Y\in {\cal B}_{1}\left( {\cal H}^{N}\right)
\end{equation}%
Bases for these two spaces can be constructed from bases of ${\cal H}^{N}$
so that 
\begin{equation}
{\cal B}_{1}\left( {\cal H}^{N}\right) ={\rm cls}\left\{ \sum\limits_{1\leq
i,j\leq \infty }\lambda _{ij}\left| \Psi _{i}\right\rangle \left\langle \Psi
_{j}\right| \right\}  \label{TCBasis}
\end{equation}%
\begin{equation}
{\cal B}\left( {\cal H}^{N}\right) ={\rm cls}\left\{ \sum_{1\leq i,j\leq
\infty }\lambda _{ij}\left| \Psi _{i}\right\rangle \left\langle \Psi
_{j}\right| \right\}  \label{BoundBasis}
\end{equation}%
where ${\rm cls\equiv closed\,linear\,span,}$ $\left\{ \Psi _{i}:1\leq i\leq
\infty \right\} $ is a complete orthonormal basis for ${\cal H}^{N}$ and the
closures are taken wrt the norms $\left\| {}\right\| _{1}$ and $\left\|
{}\right\| $ respectively. The space$\ {\cal B}_{1}\left( {\cal H}%
^{N}\right) $ forms a two-sided ideal in ${\cal B}\left( {\cal H}^{N}\right) %
\cite{ideal}$. We can also introduce the common overcomplete basis $\left\{
\left| {\sf z}\right\rangle \left\langle {\sf z}^{\prime }\right| \right\} $
formed by the generalized eigenvectors of the combined coordinate operator
(nuclear space-spin and electron space-spin). Any bounded operator can be
expanded in this basis as 
\begin{equation}
X=\int X\left( {\sf z}{\bf ,}{\sf z}^{\prime }\right) \left| {\sf z}%
\right\rangle \left\langle {\sf z}^{\prime }\right| d{\sf z}d{\sf z}^{\prime
}  \label{overcomplete}
\end{equation}%
where {\sf z}${\bf \equiv }\kappa _{1,}\ldots ,\kappa _{N_{e}},\kappa
_{N_{e}+1,}\ldots ,\kappa _{N}$, and $X\left( {\sf z}{\bf ,}{\sf z}^{\prime
}\right) $ is the integral kernel of $X$ which can be a distribution or even
more singular. As we are dealing with dual spaces the equality sign in Eq.(%
\ref{overcomplete}) needs some discussion, in this expansion it means that%
\begin{equation}
\lim_{\Delta \rightarrow {\Bbb K{\Large \times }K}}\left\| X-\int_{\Delta
}X\left( {\sf z}{\bf ,}{\sf z}^{\prime }\right) \left| {\sf z}\right\rangle
\left\langle {\sf z}^{\prime }\right| d{\sf z}d{\sf z}^{\prime }\right\| =0
\end{equation}%
where ${\Bbb K}$ $=\left\{ {\sf z}\right\} $, i.e. the set of all
coordinates. As $\left\| X\right\| \leq \left\| X\right\| _{1}$ convergence
in the trace class topology implies convergence in the operator topology,
but obviously not vice-versa, which is consistent with ${\cal B}_{1}\left( 
{\cal H}^{N}\right) \subseteq {\cal B}\left( {\cal H}^{N}\right) $. Hence if
a trace class operator is expanded in the basis $\left\{ \left| {\sf z}%
\right\rangle \left\langle {\sf z}^{\prime }\right| \right\} $ as an element
of ${\cal B}_{1}\left( {\cal H}^{N}\right) $ it has the same expansion
coefficients, i.e. integral kernel, when considered as an element of ${\cal B%
}\left( {\cal H}^{N}\right) $ i.e. as a bounded operator and its integral
kernel is absolute integrable i.e. $\int \left| X\left( {\sf z}{\bf ,}{\sf z}%
^{\prime }\right) \right| d$\underline{$\kappa $}$d$\underline{$\kappa $}$%
^{\prime }{\bf <\infty }$. Those elements of \ ${\cal B}\left( {\cal H}%
^{N}\right) $\ that do not belong to ${\cal B}_{1}\left( {\cal H}^{N}\right) 
$ have integral expansions in terms of integral kernels that are not
absolute integrable.

We can consider more general bases $\left\{ X_{\iota };1\leq \iota \leq
\infty \right\} $, $\left\{ Y_{\iota };1\leq \iota \leq \infty \right\} $ of 
${\cal B}\left( {\cal H}^{N}\right) $ and ${\cal B}_{1}\left( {\cal H}%
^{N}\right) $ than those of Eqs.(\ref{TCBasis}) and (\ref{BoundBasis}) and
expand each basis element in terms of the coordinate overcomplete set as%
\begin{eqnarray}
X_{\iota } &=&\int X_{\iota }\left( {\sf z}{\bf ,}{\sf z}^{\prime }\right)
\left| {\sf z}\right\rangle \left\langle {\sf z}^{\prime }\right| d{\sf z}d%
{\sf z}^{\prime }{\bf \equiv }\int X_{\iota }\left( {\sf z}{\bf ,}{\sf z}%
^{\prime }\right) \left( {\sf z}{\bf ,}{\sf z}^{\prime }\right| d{\sf z}d%
{\sf z}^{\prime }{\bf \equiv }\left( X_{\iota }\right|  \nonumber \\
Y_{\iota } &=&\int Y_{\iota }\left( {\sf z}{\bf ,}{\sf z}^{\prime }\right)
\left| {\sf z}\right\rangle \left\langle {\sf z}^{\prime }\right| d{\sf z}d%
{\sf z}^{\prime }{\bf \equiv }\int Y_{\iota }\left( {\sf z}{\bf ,}{\sf z}%
^{\prime }\right) \left| {\sf z}{\bf ,}{\sf z}^{\prime }\right) d{\sf z}d%
{\sf z}^{\prime }{\bf \equiv }\left| Y_{\iota }\right)  \label{Overcomplete}
\end{eqnarray}%
where we have introduced the notation $\left| {\sf z}{\bf ,}{\sf z}^{\prime
}\right) \equiv \left| {\sf z}\right\rangle \left\langle {\sf z}^{\prime
}\right| $ and the {\it super vector} notation $\left| Y\right) $ for an
operator $Y\in {\cal B}_{1}\left( {\cal H}^{N}\right) $ and $\left( {\sf z}%
{\bf ,}{\sf z}^{\prime }\right| $ and $\left( X\right| $ for operators
belonging to the dual space ${\cal B}\left( {\cal H}^{N}\right) .$ The
operators $\left\{ \left| {\sf z}\right\rangle \left\langle {\sf z}^{\prime
}\right| \right\} $ are common to both spaces and thus one can consider them
as elements of either, however the type of convergence for the continuous
linear combinations of $\left| {\sf z}\right\rangle \left\langle {\sf z}%
^{\prime }\right| $ is implied by the use of either $\left( {\sf z}{\bf ,}%
{\sf z}^{\prime }\right| $ or $\left| {\sf z}{\bf ,}{\sf z}^{\prime }\right) 
$.

We now are in a position to express maps $\hat{T}$: ${\cal B}_{1}\left( 
{\cal H}^{N}\right) \longrightarrow {\cal B}_{1}\left( {\cal H}^{N}\right) $
(sometimes called {\it super operators}) in terms of these dual bases as%
\begin{equation}
\hat{T}=\sum\Sb 1\leq \iota \leq \infty  \\ 1\leq \kappa \leq \infty  \endSb %
T_{\kappa \iota }\left| Y_{\kappa }\right) \left( X_{\iota }\right|
\end{equation}%
where the bases $\left\{ X_{\kappa }\right\} $ and $\left\{ Y_{\iota
}\right\} $ have been constructed to be biorthogonal i.e.%
\begin{equation}
\left( X_{\kappa }|Y_{\iota }\right) =\delta _{\kappa \iota }
\end{equation}%
and $T_{\kappa \iota }=\left( X_{\kappa }|\hat{T}Y_{\iota }\right) $. (A
remark on notation: we denote operators :${\cal B}_{1}\left( {\cal H}%
^{N}\right) \longrightarrow {\cal B}_{1}\left( {\cal H}^{N}\right) $ with
the decoration $\hat{T}$, while operators :${\cal H}^{N}\rightarrow {\cal H}%
^{N}$ will be denoted by the normal math font $X$ or with the decoration $%
\breve{X}$ if the domain of operation is not self evident.) An important
point to note is that unlike operators on Hilbert spaces, which are self
dual, the adjoint operation on $\hat{T}$ changes it from an operator ${\cal B%
}_{1}\left( {\cal H}^{N}\right) \longrightarrow {\cal B}_{1}\left( {\cal H}%
^{N}\right) $ to an operator ${\cal B}\left( {\cal H}^{N}\right)
\longrightarrow {\cal B}\left( {\cal H}^{N}\right) .$ Using the expansions
Eq.(\ref{Overcomplete}) one can express $\hat{T}$ as%
\begin{eqnarray}
\hat{T} &=&\sum\Sb 1\leq \iota \leq \infty  \\ 1\leq \kappa \leq \infty 
\endSb T_{\kappa \iota }\int Y_{\iota }\left( {\sf z}_{1}{\bf ,}{\sf z}%
_{2}\right) \left| {\sf z}_{1}{\bf ,}{\sf z}_{2}\right) d{\sf z}_{1}d{\sf z}%
_{2}\int X_{\kappa }^{\ast }\left( {\sf z}_{3}{\bf ,}{\sf z}_{4}\right)
\left( {\sf z}_{3}{\bf ,}{\sf z}_{4}\right| d{\sf z}_{3}d{\sf z}_{4} 
\nonumber \\
&=&\sum\Sb 1\leq \iota \leq \infty  \\ 1\leq \kappa \leq \infty  \endSb %
T_{\kappa \iota }\int Y_{\iota }\left( {\sf z}_{1}{\bf ,}{\sf z}_{2}\right)
X_{\kappa }^{\ast }\left( {\sf z}_{3}{\bf ,}{\sf z}_{4}\right) \left| {\sf z}%
_{1}{\bf ,}{\sf z}_{2}\right) \left( {\sf z}_{3}{\bf ,}{\sf z}_{4}\right| d%
{\sf z}_{1}d{\sf z}_{2}d{\sf z}_{3}d{\sf z}_{4}  \nonumber \\
&=&\int T\left( {\sf z}_{1}{\bf ,}{\sf z}_{2},{\sf z}_{3}{\bf ,}{\sf z}%
_{4}\right) \left| {\sf z}_{1}{\bf ,}{\sf z}_{2}\right) \left( {\sf z}_{3}%
{\bf ,}{\sf z}_{4}\right| d{\sf z}_{1}d{\sf z}_{2}d{\sf z}_{3}d{\sf z}_{4}
\label{SuperOp}
\end{eqnarray}%
where 
\begin{equation}
T\left( {\sf z}_{1}{\bf ,}{\sf z}_{2},{\sf z}_{3}{\bf ,}{\sf z}_{4}\right)
=\sum\Sb 1\leq \iota \leq \infty  \\ 1\leq \kappa \leq \infty  \endSb %
T_{\kappa \iota }Y_{\iota }\left( {\sf z}_{1}{\bf ,}{\sf z}_{2}\right)
X_{\kappa }^{\ast }\left( {\sf z}_{3}{\bf ,}{\sf z}_{4}\right)
\end{equation}%
and thus%
\begin{equation}
T_{\kappa \iota }=\int T\left( {\sf z}_{1}{\bf ,}{\sf z}_{2},{\sf z}_{3}{\bf %
,}{\sf z}_{4}\right) Y_{\kappa }\left( {\sf z}_{1}{\bf ,}{\sf z}_{2}\right)
X_{\iota }^{\ast }\left( {\sf z}_{3}{\bf ,}{\sf z}_{4}\right) d{\sf z}_{1}d%
{\sf z}_{2}d{\sf z}_{3}d{\sf z}_{4}
\end{equation}%
We can express Eq.(\ref{SuperOp}) in a more transparent form as%
\begin{equation}
\hat{T}=\int T\left( {\sf z}_{1}{\bf ,}{\sf z}_{2},{\sf z}_{3}{\bf ,}{\sf z}%
_{4}\right) |\left| {\sf z}_{1}\right\rangle \left\langle {\sf z}_{2}\right| 
{\bf )}\left( \left| {\sf z}_{3}\right\rangle \left\langle {\sf z}%
_{4}\right| \right| d{\sf z}_{1}d{\sf z}_{2}d{\sf z}_{3}d{\sf z}_{4}
\end{equation}%
Note that $T\left( {\sf z}_{1}{\bf ,}{\sf z}_{2},{\sf z}_{3}{\bf ,}{\sf z}%
_{4}\right) $ inherits properties from $\left\{ X_{\kappa }^{\ast }\left( 
{\sf z}_{3}{\bf ,}{\sf z}_{4}\right) \right\} $ for the $\left( {\sf z}_{3}%
{\bf ,}{\sf z}_{4}\right) $ arguments and properties from $\left\{ Y_{\iota
}\left( {\sf z}_{1}{\bf ,}{\sf z}_{2}\right) \right\} $ for the $\left( {\sf %
z}_{1}{\bf ,}{\sf z}_{2}\right) $ arguments {\it and these properties can be
quite different}. Thus these operators do not, in general, have symmetry
properties relating the $\left( {\sf z}_{1}{\bf ,}{\sf z}_{2}\right) $
arguments to the $\left( {\sf z}_{3}{\bf ,}{\sf z}_{4}\right) $ arguments.

\subsection{Operators preserving the direct sum structure ${\cal B}%
_{1}\left( {\cal H}^{N}\right) =K_{OD}\oplus K_{D}\oplus S_{D}$}

The space of trace class operators ${\cal B}_{1}\left( {\cal H}^{N}\right) $
can be decomposed into the direct sum of three subspaces defined as%
\begin{eqnarray}
K_{OD} &=&{\rm cls}\left\{ Y\in {\cal B}_{1}\left( {\cal H}^{N}\right)
:Y=\int Y\left( {\sf z}{\bf ,}{\sf z}^{\prime }\right) \left| {\sf z}%
\right\rangle \left\langle {\sf z}^{\prime }\right| d{\sf z}d{\sf z}^{\prime
}{\bf ;\quad }Y\left( {\sf z},{\sf z}^{\prime }\right) {\bf =0}\text{%
\thinspace\ if }\,{\sf z}{\bf =}{\sf z}^{\prime }\right\} \\
K_{D} &=&{\rm cls}\left\{ Y\in {\cal B}_{1}\left( {\cal H}^{N}\right)
:Y=\int Y\left( {\sf z}\right) \left| {\sf z}\right\rangle \left\langle {\sf %
z}\right| d{\sf z}\right. ;  \nonumber \\
&&\left. \qquad \qquad \quad \int Y\left( \kappa _{1,}..,\kappa
_{N_{e}},\kappa _{N_{e}+1},..\kappa _{N}\right) d\kappa _{2}\ldots d\kappa
_{N_{e}}d\kappa _{N_{e}+2}\ldots d\kappa _{N}{\bf =0}\text{ }{\bf a.e.}%
\right\} \\
S_{D} &=&{\rm cls}\left\{ Y\in {\cal B}_{1}\left( {\cal H}^{N}\right)
:Y=\int Y\left( {\sf z}\right) \left| {\sf z}\right\rangle \left\langle {\sf %
z}\right| d{\sf z}\right. \text{ };  \nonumber \\
&&\left. \qquad \qquad \quad \int Y\left( \kappa _{1,}..,\kappa
_{N_{e}},\kappa _{N_{e}+1},..\kappa _{N}\right) d\kappa _{2}\ldots d\kappa
_{N_{e}}d\kappa _{N_{e}+2}\ldots d\kappa _{N}{\bf \neq 0\,a.e.}\right\}
\end{eqnarray}%
where ${\bf a.e.\equiv }$ almost everywhere i.e. except on sets of measure
zero. We can form bases for these three subspaces of ${\cal B}_{1}\left( 
{\cal H}^{N}\right) $,%
\begin{eqnarray}
K_{OD} &=&{\rm cls}\left\{ \breve{h}_{\iota }:\breve{h}_{\iota }=\int
h_{\iota }\left( {\sf z}{\bf ,}{\sf z}^{\prime }\right) \left| {\sf z}%
\right\rangle \left\langle {\sf z}^{\prime }\right| d{\sf z}d{\sf z}^{\prime
};{\sf z}{\bf \neq }{\sf z}^{\prime }\right\}  \label{hbasis} \\
K_{D} &=&{\rm cls}\left\{ \breve{k}_{\iota }:\breve{k}_{\iota }=\int
k_{\iota }\left( {\sf z}\right) \left| {\sf z}\right\rangle \left\langle 
{\sf z}\right| d{\sf z}\right.  \nonumber \\
&&\quad \quad \quad \left. \int k_{\iota }\left( \kappa _{2,}..,\kappa
_{N_{e}},\kappa _{N_{e}+1},..\kappa _{N}\right) d\kappa _{2}\ldots d\kappa
_{N_{e}}d\kappa _{N_{e}+2}\ldots d\kappa _{N}{\bf =0}\text{ }{\bf a.e.}%
\right\}  \label{kbasis} \\
S_{D} &=&{\rm cls}\left\{ \breve{s}_{\iota }:\breve{s}_{\iota }=\int
s_{\iota }\left( {\sf z}\right) \left| {\sf z}\right\rangle \left\langle 
{\sf z}\right| d{\sf z}\right. ;  \nonumber \\
&&\quad \left. \qquad \int s_{\iota }\left( \kappa _{2,}..,\kappa
_{N_{e}},\kappa _{N_{e}+1},..\kappa _{N}\right) d\kappa _{2}\ldots d\kappa
_{N_{e}}d\kappa _{N_{e}+2}\ldots d\kappa _{N}{\bf \neq 0\,a.e.}\right\}
\label{sbasis}
\end{eqnarray}

In the following we will be considering representations in both the
overcomplete continuous basis $\left\{ \left| {\cal K}\right) \right\} $ and
the complete discrete basis $\left\{ \left| {\bf \breve{h}}\right) ,\left| 
{\bf \breve{k}}\right) ,\left| {\bf \breve{s}}\right) \right\} $ of ${\cal B}%
_{1}\left( {\cal H}^{N}\right) $, where arrays of bases vectors are denoted
by $\left| {\cal K}\right) \equiv \left| {\sf z}_{1},{\sf z}_{2}\right)
,..,\left| {\sf z}_{1}^{\prime },{\sf z}_{2}^{\prime }\right) ,..;$forming a
continuous array of uncountably many basis vectors and $\left| {\bf \breve{h}%
}\right) \equiv \left| \breve{h}_{1}\right) ,..,\left| \breve{h}_{r}\right)
;\left| {\bf \breve{k}}\right) \equiv \left| \breve{k}_{1}\right) ,..,\left| 
\breve{k}_{r}\right) $and $\left| {\bf \breve{s}}\right) \equiv \left| 
\breve{s}_{1}\right) ,..,\left| \breve{s}_{r}\right) $ forming a discrete
array of countably many basis vectors, where in general $r=\infty $. The
form of the transformation from the discrete basis to the continuous one can
be deduced by using the biorthogonal bases $\left\{ \left\{ \left( {\bf 
\breve{h}}^{d}\right| ,\left( {\bf \breve{k}}^{d}\right| ,\left( {\bf \breve{%
s}}^{d}\right| \right\} ,\left\{ \left| {\bf \breve{h}}\right) ,\left| {\bf 
\breve{k}}\right) ,\left| {\bf \breve{s}}\right) \right\} \right\} $ and $%
\left\{ \left\{ \left( {\cal K}\right| \right\} ,\left\{ \left| {\cal K}%
\right) \right\} \right\} $, (note that $\left\{ {\bf \breve{h}}^{d},{\bf 
\breve{k}}^{d},{\bf \breve{s}}^{d}\right\} \in {\cal B}\left( {\cal H}%
^{N}\right) $, while $\left\{ {\bf \breve{h}},{\bf \breve{k}},{\bf \breve{s}}%
\right\} \in {\cal B}_{1}\left( {\cal H}^{N}\right) $), and the resolutions
of the identity that they form in the following manner 
\begin{equation}
\left| Y\right) =\left| {\bf \breve{h}}\right) \left( {\bf \breve{h}}%
^{d}|Y\right) +\left| {\bf \breve{k}}\right) \left( {\bf \breve{k}}%
^{d}|Y\right) +\left| {\bf \breve{s}}\right) \left( {\bf \breve{s}}%
^{d}|Y\right) =\left| {\cal K}\right) \left( {\cal K}|Y\right)
\end{equation}%
\begin{eqnarray}
\Rightarrow \left| Y\right) &=&\left| {\cal K}\right) \left( {\cal K}|{\bf 
\breve{h}}\right) \left( {\bf \breve{h}}^{d}|Y\right) +\left| {\cal K}%
\right) \left( {\cal K}|{\bf \breve{k}}\right) \left( {\bf \breve{k}}%
^{d}|Y\right) +\left| {\cal K}\right) \left( {\cal K}|{\bf \breve{s}}\right)
\left( {\bf \breve{s}}^{d}|Y\right)  \nonumber \\
&=&\left( \left| {\cal K}\right) \right) \left( 
\begin{array}{ccc}
\left( {\cal K}|{\bf \breve{h}}\right) & \left( {\cal K}|{\bf \breve{k}}%
\right) & \left( {\cal K}|{\bf \breve{s}}\right)%
\end{array}%
\right) \left( 
\begin{array}{c}
\left( {\bf \breve{h}}^{d}|Y\right) \\ 
\left( {\bf \breve{k}}^{d}|Y\right) \\ 
\left( {\bf \breve{s}}^{d}|Y\right)%
\end{array}%
\right)  \label{KtoH}
\end{eqnarray}%
and 
\begin{eqnarray}
\Rightarrow \left| Y\right) &=&\left| {\bf \breve{h}}\right) \left( {\bf 
\breve{h}}^{d}|{\cal K}\right) \left( {\cal K}|Y\right) +\left| {\bf \breve{k%
}}\right) \left( {\bf \breve{k}}^{d}|{\cal K}\right) \left( {\cal K}%
|Y\right) +\left| {\bf \breve{s}}\right) \left( {\bf \breve{s}}^{d}|{\cal K}%
\right) \left( {\cal K}|Y\right)  \nonumber \\
&=&\left( \left| {\bf \breve{h}}\right) ,\left| {\bf \breve{k}}\right)
,\left| {\bf \breve{s}}\right) \right) \left( 
\begin{array}{c}
\left( {\bf \breve{h}}^{d}|{\cal K}\right) \\ 
\left( {\bf \breve{k}}^{d}|{\cal K}\right) \\ 
\left( {\bf \breve{s}}^{d}|{\cal K}\right)%
\end{array}%
\right) \left( {\cal K}|Y\right)  \label{HtoK}
\end{eqnarray}%
where 
\begin{eqnarray}
\left| {\bf \breve{h}}\right) \left( {\bf \breve{h}}^{d}\right|
&=&\sum_{1\leq \iota \leq \infty }\left| \breve{h}_{\iota }\right) \left( 
\breve{h}_{\iota }^{d}\right| =\hat{P}_{K_{OD}}  \nonumber \\
\left| {\bf \breve{k}}\right) \left( {\bf \breve{k}}^{d}\right|
&=&\sum_{1\leq \iota \leq \infty }\left| \breve{k}_{\iota }\right) \left( 
\breve{k}_{\iota }^{d}\right| =\hat{P}_{K_{D}}  \nonumber \\
\left| {\bf \breve{s}}\right) \left( {\bf \breve{s}}^{d}\right|
&=&\sum_{1\leq \iota \leq \infty }\left| \breve{s}_{\iota }\right) \left( 
\breve{s}_{\iota }^{d}\right| =\hat{P}_{S_{D}}  \nonumber \\
I &=&\left| {\bf \breve{h}}\right) \left( {\bf \breve{h}}^{d}\right| +\left| 
{\bf \breve{k}}\right) \left( {\bf \breve{k}}^{d}\right| +\left| {\bf \breve{%
s}}\right) \left( {\bf \breve{s}}^{d}\right|  \nonumber \\
I &=&\left| {\cal K}\right) \left( {\cal K}\right| =\int \left| {\sf z}_{1},%
{\sf z}_{2}\right) \left( {\sf z}_{1},{\sf z}_{2}\right| d{\sf z}_{1}d{\sf z}%
_{2}
\end{eqnarray}%
$\hat{P}_{K_{OD}}$, $\hat{P}_{K_{D}}$ and $\hat{P}_{S_{D}}$ are the
projectors onto $K_{OD}$, $K_{D}$ and $S_{D}$ respectively and their
integral kernels are given by 
\begin{eqnarray}
P_{K_{OD}}\left( {\sf z}_{1},{\sf z}_{2},{\sf z}_{3},{\sf z}_{4}\right)
&=&\sum_{\kappa }\left( {\sf z}_{1},{\sf z}_{2}|\breve{h}_{\kappa }\right)
\left( \breve{h}_{\kappa }^{d}|{\sf z}_{3},{\sf z}_{4}\right) \\
P_{K_{D}}\left( {\sf z}_{1},{\sf z}_{2},{\sf z}_{3},{\sf z}_{4}\right)
&=&\sum_{\kappa }\left( {\sf z}_{1},{\sf z}_{2}|\breve{k}_{\kappa }\right)
\left( \breve{k}_{\kappa }^{d}|{\sf z}_{3},{\sf z}_{4}\right)  \nonumber \\
&=&\sum_{\kappa }\left( {\sf z}_{1},{\sf z}_{1}|\breve{k}_{\kappa }\right)
\left( \breve{k}_{\kappa }^{d}|{\sf z}_{3},{\sf z}_{3}\right) \delta \left( 
{\sf z}_{1}-{\sf z}_{2}\right) \delta \left( {\sf z}_{3}-{\sf z}_{4}\right) 
\nonumber \\
&\equiv &P_{K_{D}}\left( {\sf z}_{1},{\sf z}_{3}\right) \delta \left( {\sf z}%
_{1}-{\sf z}_{2}\right) \delta \left( {\sf z}_{3}-{\sf z}_{4}\right)
\end{eqnarray}%
\begin{eqnarray}
P_{S_{D}}\left( {\sf z}_{1},{\sf z}_{2},{\sf z}_{3},{\sf z}_{4}\right)
&=&\sum_{\kappa }\left( {\sf z}_{1},{\sf z}_{2}|\breve{s}_{\kappa }\right)
\left( \breve{s}_{\kappa }^{d}|{\sf z}_{3},{\sf z}_{4}\right)  \nonumber \\
&=&\sum_{\kappa }\left( {\sf z}_{1},{\sf z}_{1}|\breve{s}_{\kappa }\right)
\left( \breve{s}_{\kappa }^{d}|{\sf z}_{3},{\sf z}_{3}\right) \delta \left( 
{\sf z}_{1}-{\sf z}_{2}\right) \delta \left( {\sf z}_{3}-{\sf z}_{4}\right) 
\nonumber \\
&\equiv &P_{S_{D}}\left( {\sf z}_{1},{\sf z}_{3}\right) \delta \left( {\sf z}%
_{1}-{\sf z}_{2}\right) \delta \left( {\sf z}_{3}-{\sf z}_{4}\right)
\end{eqnarray}%
Note that these projectors are {\bf not }orthogonal projectors and {\it in
fact }there is no property of orthogonality nor adjoint operation on
operators acting in ${\cal B}_{1}\left( {\cal H}^{N}\right) $, as it is{\it %
\ not }a Hilbert space, unless ${\cal H}^{N}$ is finite dimensional. They do
however have the properties 
\begin{eqnarray}
\hat{P}_{A}^{2} &=&\hat{P}_{A}  \nonumber \\
\hat{P}_{A}\hat{P}_{B} &=&0\quad A\neq B
\end{eqnarray}%
The matrix transforming from the $\left\{ \left| {\bf \breve{h}}\right)
,\left| {\bf \breve{k}}\right) ,\left| {\bf \breve{s}}\right) \right\} $
basis to the $\left\{ \left| {\cal K}\right) \right\} $ basis obtained from
Eq.(\ref{KtoH}) is, 
\begin{equation}
{\Bbb S=}\left( 
\begin{array}{c}
\left( {\bf \breve{h}}^{d}|{\cal K}\right) \\ 
\left( {\bf \breve{k}}^{d}|{\cal K}\right) \\ 
\left( {\bf \breve{s}}^{d}|{\cal K}\right)%
\end{array}%
\right)
\end{equation}%
i.e.%
\begin{equation}
\left| {\cal K}\right) =\left( \left| {\bf \breve{h}}\right) ,\left| {\bf 
\breve{k}}\right) ,\left| {\bf \breve{s}}\right) \right) \left( 
\begin{array}{c}
\left( {\bf \breve{h}}^{d}|{\cal K}\right) \\ 
\left( {\bf \breve{k}}^{d}|{\cal K}\right) \\ 
\left( {\bf \breve{s}}^{d}|{\cal K}\right)%
\end{array}%
\right)
\end{equation}%
and using Eq.(\ref{HtoK}) the transformation from the $\left\{ \left| {\cal K%
}\right) \right\} $ to the $\left\{ \left| {\bf \breve{h}}\right) ,\left| 
{\bf \breve{k}}\right) ,\left| {\bf \breve{s}}\right) \right\} $ basis is 
\begin{equation}
{\Bbb S}^{-1}{\Bbb =}\left( 
\begin{array}{ccc}
\left( {\cal K}|{\bf \breve{h}}\right) & \left( {\cal K}|{\bf \breve{k}}%
\right) & \left( {\cal K}|{\bf \breve{s}}\right)%
\end{array}%
\right)
\end{equation}%
i.e.%
\begin{equation}
\left( \left| {\bf \breve{h}}\right) ,\left| {\bf \breve{k}}\right) ,\left| 
{\bf \breve{s}}\right) \right) =\left| {\cal K}\right) \left( 
\begin{array}{ccc}
\left( {\cal K}|{\bf \breve{h}}\right) & \left( {\cal K}|{\bf \breve{k}}%
\right) & \left( {\cal K}|{\bf \breve{s}}\right)%
\end{array}%
\right)
\end{equation}%
{\bf Note }that ${\Bbb S}^{-1}$ {\it is not }the transpose or adjoint matrix
of ${\Bbb S}$ as ${\cal B}_{1}\left( {\cal H}^{N}\right) $ is not self dual,
one can for instance can see that ${\Bbb S}^{-1}$ involves integral kernels $%
\left( {\sf z}_{1},{\sf z}_{2}|\breve{h}_{\iota }\right) $ of trace class
operators, while ${\Bbb S}$ involves integral kernels $\left( \breve{h}%
_{\iota }^{d}|{\sf z}_{1},{\sf z}_{2}\right) $ of bounded operators that do
not have to be of trace class. We also have the biorthogonality condition 
\begin{eqnarray}
&&\left( 
\begin{array}{c}
\left( {\bf \breve{h}}^{d}\right| \\ 
\left( {\bf \breve{k}}^{d}\right| \\ 
\left( {\bf \breve{s}}^{d}\right|%
\end{array}%
\right) \left( \left| {\bf \breve{h}}\right) ,\left| {\bf \breve{k}}\right)
,\left| {\bf \breve{s}}\right) \right) =\left( 
\begin{array}{ccc}
\left( {\bf \breve{h}}^{d}|{\bf \breve{h}}\right) & \left( {\bf \breve{h}}%
^{d}|{\bf \breve{k}}\right) & \left( {\bf \breve{h}}^{d}|{\bf \breve{s}}%
\right) \\ 
\left( {\bf \breve{k}}^{d}|{\bf \breve{h}}\right) & \left( {\bf \breve{k}}%
^{d}|{\bf \breve{k}}\right) & \left( {\bf \breve{k}}^{d}|{\bf \breve{s}}%
\right) \\ 
\left( {\bf \breve{s}}^{d}|{\bf \breve{h}}\right) & \left( {\bf \breve{s}}%
^{d}|{\bf \breve{k}}\right) & \left( {\bf \breve{s}}^{d}|{\bf \breve{s}}%
\right)%
\end{array}%
\right) ={\Bbb I} \\
&\Leftrightarrow &{\Bbb S}^{d}\left( 
\begin{array}{c}
\left( {\bf \breve{h}}^{d}\right| \\ 
\left( {\bf \breve{k}}^{d}\right| \\ 
\left( {\bf \breve{s}}^{d}\right|%
\end{array}%
\right) \left( \left| {\bf \breve{h}}\right) ,\left| {\bf \breve{k}}\right)
,\left| {\bf \breve{s}}\right) \right) {\Bbb S}=\left( {\cal K}|{\cal K}%
\right) ={\cal I=}{\Bbb S}^{d}{\Bbb S}
\end{eqnarray}%
where ${\cal I}$ is the integral kernel $\delta \left( {\sf z}_{1}-{\sf z}%
_{3}\right) \delta \left( {\sf z}_{2}-{\sf z}_{4}\right) $ and ${\Bbb S}%
^{d}: $ ${\cal B}\left( {\cal H}^{N}\right) \rightarrow {\cal B}\left( {\cal %
H}^{N}\right) $ is the transformation between the dual bases $\left( 
\begin{array}{c}
\left( {\bf \breve{h}}^{d}\right| \\ 
\left( {\bf \breve{k}}^{d}\right| \\ 
\left( {\bf \breve{s}}^{d}\right|%
\end{array}%
\right) \rightarrow \left( {\cal K}\right| $, which shows that ${\Bbb S}^{d}=%
{\Bbb S}^{-1}$. Hence we obtain the identities 
\begin{eqnarray}
&&{\cal \hat{M}}\breve{Y}=\left( \left| {\bf \breve{h}}\right) ,\left| {\bf 
\breve{k}}\right) ,\left| {\bf \breve{s}}\right) \right) {\Bbb SS}%
^{-1}\left( 
\begin{array}{ccc}
{\Bbb M}_{K_{OD}} & {\Bbb M}_{K_{OD}K_{D}} & {\Bbb M}_{K_{OD}S_{D}} \\ 
{\Bbb M}_{K_{D}K_{OD}} & {\Bbb M}_{K_{D}} & {\Bbb M}_{K_{D}S_{D}} \\ 
{\Bbb M}_{S_{D}K_{OD}} & {\Bbb M}_{S_{D}K_{D}} & {\Bbb T}_{S_{D}}%
\end{array}%
\right) {\Bbb SS}^{-1}\left( 
\begin{array}{c}
{\Bbb Y}_{_{K_{OD}}} \\ 
{\Bbb Y}_{_{K_{D}}} \\ 
{\Bbb Y}_{_{S_{D}}}%
\end{array}%
\right) = \\
&&  \nonumber \\
&&\left( \left| {\cal K}\right) \right) {\Bbb S}^{-1}\left( 
\begin{array}{ccc}
{\Bbb M}_{K_{OD}} & {\Bbb M}_{K_{OD}K_{D}} & {\Bbb M}_{K_{OD}S_{D}} \\ 
{\Bbb M}_{K_{D}K_{OD}} & {\Bbb M}_{K_{D}} & {\Bbb M}_{K_{D}S_{D}} \\ 
{\Bbb M}_{S_{D}K_{OD}} & {\Bbb M}_{S_{D}K_{D}} & {\Bbb T}_{S_{D}}%
\end{array}%
\right) {\Bbb S}\left( {\cal K}|Y\right) =\left( \left| {\cal K}\right)
\right) \left( {\cal K}|{\cal \hat{M}K}\right) \left( {\cal K}|Y\right)
\end{eqnarray}%
where 
\begin{equation}
\left( {\cal K}|{\cal \hat{M}K}\right) =\left( 
\begin{array}{ccc}
\left( {\cal K}|{\bf \breve{h}}\right) & \left( {\cal K}|{\bf \breve{k}}%
\right) & \left( {\cal K}|{\bf \breve{s}}\right)%
\end{array}%
\right) \left( 
\begin{array}{ccc}
{\Bbb M}_{K_{OD}} & {\Bbb M}_{K_{OD}K_{D}} & {\Bbb M}_{K_{OD}S_{D}} \\ 
{\Bbb M}_{K_{D}K_{OD}} & {\Bbb M}_{K_{D}} & {\Bbb M}_{K_{D}S_{D}} \\ 
{\Bbb M}_{S_{D}K_{OD}} & {\Bbb M}_{S_{D}K_{D}} & {\Bbb M}_{S_{D}}%
\end{array}%
\right) \left( 
\begin{array}{c}
\left( {\bf \breve{h}}^{d}|{\cal K}\right) \\ 
\left( {\bf \breve{k}}^{d}|{\cal K}\right) \\ 
\left( {\bf \breve{s}}^{d}|{\cal K}\right)%
\end{array}%
\right)
\end{equation}%
We also have that 
\begin{eqnarray}
&&{\cal \hat{M}}\breve{Y}=\left( \left| {\cal K}\right) \right) {\Bbb S}%
^{-1}\left( 
\begin{array}{ccc}
{\Bbb M}_{K_{OD}} & {\Bbb M}_{K_{OD}K_{D}} & {\Bbb M}_{K_{OD}S_{D}} \\ 
{\Bbb M}_{K_{D}K_{OD}} & {\Bbb M}_{K_{D}} & {\Bbb M}_{K_{D}S_{D}} \\ 
{\Bbb M}_{S_{D}K_{OD}} & {\Bbb M}_{S_{D}K_{D}} & {\Bbb M}_{S_{D}}%
\end{array}%
\right) {\Bbb SS}^{-1}\left( 
\begin{array}{c}
{\Bbb Y}_{_{K_{OD}}} \\ 
{\Bbb Y}_{_{K_{D}}} \\ 
{\Bbb Y}_{_{S_{D}}}%
\end{array}%
\right) =  \nonumber \\
&&  \nonumber \\
&&\left( \left| {\cal K}\right) \right) \left\{ \left( {\cal K}|{\cal \hat{M}%
}_{K_{OD}}{\cal K}\right) +\left( {\cal K}|{\cal \hat{M}}_{K_{OD}K_{D}}{\cal %
K}\right) +\left( {\cal K}|{\cal \hat{M}}_{K_{OD}S_{D}}{\cal K}\right)
+\right.  \nonumber \\
&&\qquad \left( {\cal K}|{\cal \hat{M}}_{K_{D}K_{OD}}{\cal K}\right) +\left( 
{\cal K}|{\cal \hat{M}}_{K_{D}}{\cal K}\right) +\left( {\cal K}|{\cal \hat{M}%
}_{K_{D}S_{D}}{\cal K}\right) +\left( {\cal K}|{\cal \hat{M}}_{S_{D}K_{OD}}%
{\cal K}\right) +  \nonumber \\
&&\qquad \qquad \left. \left( {\cal K}|{\cal \hat{M}}_{S_{D}K_{D}}{\cal K}%
\right) +\left( {\cal K}|{\cal \hat{M}}_{S_{D}}{\cal K}\right) \right\}
\left\{ \left( {\cal K}|Y_{K_{OD}}\right) +\left( {\cal K}|Y_{K_{D}}\right)
+\left( {\cal K}|Y_{S_{D}}\right) \right\}  \label{func}
\end{eqnarray}%
where 
\begin{equation}
{\cal \hat{M}}_{AB}=\hat{P}_{A}{\cal \hat{M}}\hat{P}_{B}\quad \text{and\quad 
}Y_{A}=\hat{P}_{A}Y
\end{equation}%
for various subspaces $A$ and $B$ of ${\cal B}_{1}\left( {\cal H}^{N}\right) 
$. The expression in Eq.(\ref{func}) can be further developed to give 
\begin{equation}
{\cal \hat{M}}\breve{Y}=\breve{Y}^{\prime }=\left( \left| {\cal K}\right)
\right) \left\{ \left( {\cal K}|Y_{K_{OD}}^{\prime }\right) +\left( {\cal K}%
|Y_{K_{D}}^{\prime }\right) +\left( {\cal K}|Y_{S_{D}}^{\prime }\right)
\right\}
\end{equation}%
where 
\begin{eqnarray}
\left( 
\begin{array}{c}
\left( {\cal K}|Y_{K_{OD}}^{\prime }\right) \\ 
\left( {\cal K}|Y_{K_{D}}^{\prime }\right) \\ 
\left( {\cal K}|Y_{S_{D}}^{\prime }\right)%
\end{array}%
\right) &=&\left( 
\begin{array}{ccc}
\left( {\cal K}|{\cal \hat{M}}_{K_{OD}}{\cal K}\right) & \left( {\cal K}|%
{\cal \hat{M}}_{K_{OD}K_{D}}{\cal K}\right) & \left( {\cal K}|{\cal \hat{M}}%
_{K_{OD}S_{D}}{\cal K}\right) \\ 
\left( {\cal K}|{\cal \hat{M}}_{K_{D}K_{OD}}{\cal K}\right) & \left( {\cal K}%
|{\cal \hat{M}}_{K_{D}}{\cal K}\right) & \left( {\cal K}|{\cal \hat{M}}%
_{S_{D}K_{OD}}{\cal K}\right) \\ 
\left( {\cal K}|{\cal \hat{M}}_{S_{D}K_{D}}{\cal K}\right) & \left( {\cal K}|%
{\cal \hat{M}}_{S_{D}K_{D}}{\cal K}\right) & \left( {\cal K}|{\cal \hat{M}}%
_{S_{D}}{\cal K}\right)%
\end{array}%
\right) \left( 
\begin{array}{c}
\left( {\cal K}|Y_{K_{OD}}\right) \\ 
\left( {\cal K}|Y_{K_{D}}\right) \\ 
\left( {\cal K}|Y_{S_{D}}\right)%
\end{array}%
\right)  \nonumber \\
&&
\end{eqnarray}

Transformations of ${\cal B}_{1}\left( {\cal H}^{N}\right) $ that preserve
the vector bundle structure of ${\frak B}_{vect}$ can be expressed as a sum%
\begin{eqnarray}
{\frak \hat{T}} &{\frak =}&{\frak \hat{T}}_{K_{OD}}+{\frak \hat{T}}_{K_{D}}+%
{\frak \hat{T}}_{S_{D}}+{\frak \hat{T}}_{K_{D}S_{D}}+{\frak \hat{T}}%
_{K_{OD}S_{D}}+{\frak \hat{T}}_{K_{OD}K_{D}}+{\frak \hat{T}}_{K_{D}K_{OD}} 
\nonumber \\
{\frak \hat{T}} &:&{\frak B}_{vect}\longrightarrow {\frak B}_{vect}
\end{eqnarray}%
where ${\frak \hat{T}}_{AB}$ denotes maps between subspaces $A$ and $B$.
These can be represented as matrices wrt the bases introduced in Eqs. (\ref%
{hbasis}),(\ref{kbasis}) and (\ref{sbasis}) as 
\begin{equation}
\left( 
\begin{array}{ccc}
{\Bbb T}_{K_{OD}} & {\Bbb T}_{K_{OD}K_{D}} & {\Bbb T}_{K_{OD}S_{D}} \\ 
{\Bbb T}_{K_{D}K_{OD}} & {\Bbb T}_{K_{D}} & {\Bbb T}_{K_{D}S_{D}} \\ 
{\Bbb O} & {\Bbb O} & {\Bbb T}_{S_{D}}%
\end{array}%
\right)
\end{equation}%
The non zero components can be expanded as%
\begin{eqnarray}
{\frak \hat{T}}_{K_{OD}} &=&\sum_{\kappa ,\iota }{\frak T}_{K_{OD}\kappa
\iota }\left| \breve{h}_{\kappa }\right) \left( \breve{h}_{\iota }^{d}\right|
\nonumber \\
{\frak \hat{T}}_{K_{D}} &=&\sum_{\kappa ,\iota }{\frak T}_{K_{D}\kappa \iota
}\left| \breve{k}_{\kappa }\right) \left( \breve{k}_{\iota }^{d}\right| 
\nonumber \\
{\frak \hat{T}}_{S_{D}} &=&\sum_{\kappa ,\iota }{\frak T}_{S_{D}\kappa \iota
}\left| \breve{s}_{\kappa }\right) \left( \breve{s}_{\iota }^{d}\right| 
\nonumber \\
{\frak \hat{T}}_{K_{D}S_{D}} &=&\sum_{\kappa ,\iota }{\frak T}%
_{K_{D}S_{D}\kappa \iota }\left| \breve{k}_{\kappa }\right) \left( \breve{s}%
_{\iota }^{d}\right|  \nonumber \\
{\frak \hat{T}}_{K_{OD}K_{D}} &=&\sum_{\kappa ,\iota }{\frak T}%
_{K_{OD}K_{D}\kappa \iota }\left| \breve{h}_{\kappa }\right) \left( \breve{k}%
_{\iota }^{d}\right|  \nonumber \\
{\frak \hat{T}}_{K_{D}K_{OD}} &=&\sum_{\kappa ,\iota }{\frak T}%
_{K_{D}K_{OD}\kappa \iota }\left| \breve{k}_{\kappa }\right) \left( \breve{h}%
_{\iota }^{d}\right|  \nonumber \\
{\frak \hat{T}}_{K_{OD}S_{D}} &=&\sum_{\kappa ,\iota }{\frak T}%
_{K_{OD}S_{D}\kappa \iota }\left| \breve{h}_{\kappa }\right) \left( \breve{s}%
_{\iota }^{d}\right|
\end{eqnarray}

In order to analyze the properties of Lie groups formed from such operators
it is convenient to express them as ${\frak \hat{T}}=\hat{I}+\hat{\Lambda}%
=e^{\hat{\Upsilon}}$, then%
\begin{eqnarray}
{\frak \hat{T}}_{K_{OD}} &=&\hat{P}_{K_{OD}}+\hat{\Lambda}_{K_{OD}}=\hat{P}%
_{K_{OD}}{\frak \hat{T}}\hat{P}_{K_{OD}}  \nonumber \\
{\frak \hat{T}}_{K_{D}} &=&\hat{P}_{K_{D}}+\hat{\Lambda}_{K_{D}}=\hat{P}%
_{K_{D}}{\frak \hat{T}}\hat{P}_{K_{D}}  \nonumber \\
{\frak \hat{T}}_{S_{D}} &=&\hat{P}_{S_{D}}+\hat{\Lambda}_{S_{D}}=\hat{P}%
_{S_{D}}{\frak \hat{T}}\hat{P}_{S_{D}}  \nonumber \\
{\frak \hat{T}}_{K_{D}S_{D}} &=&\hat{\Lambda}_{K_{D}S_{D}}=\hat{P}_{K_{D}}%
{\frak \hat{T}}\hat{P}_{S_{D}}  \nonumber \\
{\frak \hat{T}}_{K_{D}K_{OD}} &=&\hat{\Lambda}_{K_{D}K_{OD}}=\hat{P}_{K_{D}}%
{\frak \hat{T}}\hat{P}_{K_{OD}}  \nonumber \\
{\frak \hat{T}}_{K_{OD}K_{D}} &=&\hat{\Lambda}_{K_{OD}K_{D}}=\hat{P}_{K_{OD}}%
{\frak \hat{T}}\hat{P}_{K_{D}}  \nonumber \\
{\frak \hat{T}}_{K_{OD}S_{D}} &=&\hat{\Lambda}_{K_{OD}S_{D}}=\hat{P}_{K_{OD}}%
{\frak \hat{T}}\hat{P}_{S_{D}}
\end{eqnarray}%
(A notational note: in kernels where arguments are repeated we only display
one argument if this does not detract from the clarity of the exposition).
The integral kernels of ${\frak \hat{T}}_{K_{OD}}$, ${\frak \hat{T}}_{K_{D}}$%
, and ${\frak \hat{T}}_{S_{D}}$, are given by 
\begin{eqnarray}
{\frak T}_{K_{OD}}\left( {\sf z}_{1},{\sf z}_{2},{\sf z}_{3},{\sf z}%
_{4}\right) &=&\left( {\sf z}_{1},{\sf z}_{2}|{\frak \hat{T}}_{K_{OD}}{\sf z}%
_{3},{\sf z}_{4}\right) =\left( \left| {\sf z}_{1}\right\rangle \left\langle 
{\sf z}_{2}\right| |{\frak \hat{T}}_{K_{OD}}\left| {\sf z}_{3}\right\rangle
\left\langle {\sf z}_{4}\right| \right)  \nonumber \\
&=&\sum_{\kappa ,\iota }{\frak T}_{K_{OD}\kappa \iota }\left( \left| {\sf z}%
_{1}\right\rangle \left\langle {\sf z}_{2}\right| |\breve{h}_{\kappa
}\right) \left( \breve{h}_{\iota }^{d}|\left| {\sf z}_{3}\right\rangle
\left\langle {\sf z}_{4}\right| \right)  \nonumber \\
&=&\sum_{\kappa ,\iota }{\frak T}_{K_{OD}\kappa \iota }\breve{h}_{\kappa
}\left( {\sf z}_{1},{\sf z}_{2}\right) \breve{h}_{\iota }^{d\ast }\left( 
{\sf z}_{3},{\sf z}_{4}\right) \\
&=&P_{K_{OD}}\left( {\sf z}_{1},{\sf z}_{2},{\sf z}_{3},{\sf z}_{4}\right)
+\sum_{\kappa ,\iota }\Lambda _{K_{OD}\kappa \iota }\breve{h}_{\kappa
}\left( {\sf z}_{1},{\sf z}_{2}\right) \breve{h}_{\iota }^{d\ast }\left( 
{\sf z}_{3},{\sf z}_{4}\right)  \nonumber \\
&=&P_{K_{OD}}\left( {\sf z}_{1},{\sf z}_{2},{\sf z}_{3},{\sf z}_{4}\right)
+\Lambda _{K_{OD}}\left( {\sf z}_{1},{\sf z}_{2},{\sf z}_{3},{\sf z}%
_{4}\right)  \label{KOD}
\end{eqnarray}%
\begin{equation}
\qquad \qquad {\frak T}_{K_{D}}\left( {\sf z}_{1},{\sf z}_{2}\right)
=\sum_{\kappa ,\iota }{\frak T}_{K_{D}\kappa \iota }k_{\kappa }\left( {\sf z}%
_{1}\right) k_{\iota }^{d\ast }\left( {\sf z}_{2}\right) =P_{K_{D}}\left( 
{\sf z}_{1},{\sf z}_{2}\right) +\Lambda _{K_{D}}\left( {\sf z}_{1},{\sf z}%
_{2}\right)  \label{KD}
\end{equation}%
\begin{equation}
\qquad \qquad {\frak T}_{S_{D}}\left( {\sf z}_{1},{\sf z}_{2}\right)
=\sum_{\kappa ,\iota }{\frak T}_{S_{D}\kappa \iota }s_{\kappa }\left( {\sf z}%
_{1}\right) s_{\iota }^{d\ast }\left( {\sf z}_{2}\right) =P_{S_{D}}\left( 
{\sf z}_{1},{\sf z}_{2}\right) +\Lambda _{S_{D}}\left( {\sf z}_{1},{\sf z}%
_{2}\right)  \label{SD}
\end{equation}%
and those of ${\frak \hat{T}}_{K_{OD}K_{D}}$, ${\frak \hat{T}}_{K_{D}K_{OD}}$%
, ${\frak \hat{T}}_{K_{OD}S_{D}}$ and ${\frak \hat{T}}_{K_{D}S_{D}}$ by%
\begin{eqnarray}
{\frak T}_{K_{D}S_{D}}\left( {\sf z}_{1},{\sf z}_{2}\right) &=&\sum_{\kappa
,\iota }{\frak T}_{K_{D}S_{D}\kappa \iota }k_{\kappa }\left( {\sf z}%
_{1}\right) s_{\iota }^{d\ast }\left( {\sf z}_{2}\right)  \nonumber \\
&=&\sum_{\kappa ,\iota }\Lambda _{K_{D}S_{D}\kappa \iota }k_{\kappa }\left( 
{\sf z}_{1}\right) s_{\iota }^{d\ast }\left( {\sf z}_{2}\right) =\Lambda
_{K_{D}S_{D}}\left( {\sf z}_{1},{\sf z}_{2}\right)  \nonumber \\
{\frak T}_{K_{OD}S_{D}}\left( {\sf z}_{1},{\sf z}_{2},{\sf z}_{3}\right)
&=&\sum_{\kappa ,\iota }{\frak T}_{K_{OD}S_{D}\kappa \iota }h_{\kappa
}\left( {\sf z}_{1},{\sf z}_{2}\right) s_{\iota }^{d\ast }\left( {\sf z}%
_{3}\right)  \nonumber \\
&=&\sum_{\kappa ,\iota }\Lambda _{K_{OD}S_{D}\kappa \iota }h_{\kappa }\left( 
{\sf z}_{1},{\sf z}_{2}\right) s_{\iota }^{d\ast }\left( {\sf z}_{3}\right)
=\Lambda _{K_{OD}S_{D}}\left( {\sf z}_{1},{\sf z}_{2},{\sf z}_{3}\right) 
\nonumber \\
{\frak T}_{K_{D}K_{OD}}\left( {\sf z}_{1},{\sf z}_{2},{\sf z}_{3}\right)
&=&\sum_{\kappa ,\iota }{\frak T}_{K_{D}K_{OD}\kappa \iota }k_{\kappa
}\left( {\sf z}_{1}\right) h_{\iota }^{d\ast }\left( {\sf z}_{1},{\sf z}%
_{2}\right)  \nonumber \\
&=&\sum_{\kappa ,\iota }\Lambda _{K_{D}K_{OD}\kappa \iota }k_{\kappa }\left( 
{\sf z}_{1}\right) h_{\iota }^{d\ast }\left( {\sf z}_{2},{\sf z}_{3}\right)
=\Lambda _{K_{D}K_{OD}}\left( {\sf z}_{1},{\sf z}_{2},{\sf z}_{3}\right) 
\nonumber \\
{\frak T}_{K_{OD}K_{D}}\left( {\sf z}_{1},{\sf z}_{2},{\sf z}_{3}\right)
&=&\sum_{\kappa ,\iota }{\frak T}_{K_{OD}K_{D}\kappa \iota }h_{\kappa
}\left( {\sf z}_{1},{\sf z}_{2}\right) k_{\iota }^{d\ast }\left( {\sf z}%
_{3}\right)  \nonumber \\
&=&\sum_{\kappa ,\iota }\Lambda _{K_{OD}K_{D}\kappa \iota }h_{\kappa }\left( 
{\sf z}_{1},{\sf z}_{2}\right) k_{\iota }^{d\ast }\left( {\sf z}_{3}\right)
=\Lambda _{K_{OD}K_{D}}\left( {\sf z}_{1},{\sf z}_{2},{\sf z}_{3}\right)
\label{mixed}
\end{eqnarray}%
The super operator transformations, ${\frak \hat{T}}$, can be associated
with matrices of integral kernels as 
\begin{eqnarray}
{\frak \hat{T}\equiv }\left( 
\begin{array}{ccc}
{\frak T}_{K_{OD}} & {\frak T}_{K_{OD,}K_{D}} & {\frak T}_{K_{D},S_{D}} \\ 
{\frak T}_{K_{D},K_{OD}} & {\frak T}_{K_{D}} & {\frak T}_{K_{D}S_{D}} \\ 
0 & 0 & {\frak T}_{S_{D}}%
\end{array}%
\right) = &&\left( 
\begin{array}{ccc}
P_{K_{OD}} & 0 & 0 \\ 
0 & P_{K_{D}} & 0 \\ 
0 & 0 & P_{S_{D}}%
\end{array}%
\right)  \nonumber \\
&&+\left( 
\begin{array}{ccc}
\Lambda _{K_{OD}} & \Lambda _{K_{D}K_{OD}} & \Lambda _{K_{D}S_{D}} \\ 
\Lambda _{K_{D}K_{OD}} & \Lambda _{K_{D}} & \Lambda _{K_{OD}S_{D}} \\ 
0 & 0 & \Lambda _{S_{D}}%
\end{array}%
\right)
\end{eqnarray}%
which produce the transformed integral kernels $Y^{\prime }\left( {\sf z}%
_{1},{\sf z}_{2}\right) ={\frak \hat{T}}Y\left( {\sf z}_{1},{\sf z}%
_{2}\right) $ of trace class operators in the following manner%
\begin{equation}
\left( {\sf z}_{1},{\sf z}_{2}|{\frak \hat{T}}Y\right) =\left( {\sf z}_{1},%
{\sf z}_{2}|\breve{Y}_{K_{OD}}^{\prime }\right) +\left( {\sf z}_{1},{\sf z}%
_{2}|\breve{Y}_{K_{D}}^{\prime }\right) +\left( {\sf z}_{1},{\sf z}_{2}|%
\breve{Y}_{S_{D}}^{\prime }\right)  \label{transkern}
\end{equation}%
where%
\begin{eqnarray}
&&\left( 
\begin{array}{c}
\left( {\sf z}_{1},{\sf z}_{2}|\breve{Y}_{K_{OD}}^{\prime }\right) \\ 
\left( {\sf z}_{1},{\sf z}_{2}|\breve{Y}_{K_{D}}^{\prime }\right) \\ 
\left( {\sf z}_{1},{\sf z}_{2}|\breve{Y}_{S_{D}}^{\prime }\right)%
\end{array}%
\right) =  \nonumber \\
&&\int \left( 
\begin{array}{ccc}
{\frak T}_{K_{OD}}\left( {\sf z}_{1},{\sf z}_{2},{\sf z}_{3},{\sf z}%
_{4}\right) & {\frak T}_{K_{OD,}K_{D}}\left( {\sf z}_{1},{\sf z}_{2},{\sf z}%
_{3},{\sf z}_{4}\right) & {\frak T}_{K_{D},S_{D}}\left( {\sf z}_{1},{\sf z}%
_{2},{\sf z}_{3},{\sf z}_{4}\right) \\ 
{\frak T}_{K_{D},K_{OD}}\left( {\sf z}_{1},{\sf z}_{2},{\sf z}_{3},{\sf z}%
_{4}\right) & {\frak T}_{K_{D}}\left( {\sf z}_{1},{\sf z}_{2},{\sf z}_{3},%
{\sf z}_{4}\right) & {\frak T}_{K_{D}S_{D}}\left( {\sf z}_{1},{\sf z}_{2},%
{\sf z}_{3},{\sf z}_{4}\right) \\ 
0 & 0 & {\frak T}_{S_{D}}\left( {\sf z}_{1},{\sf z}_{2},{\sf z}_{3},{\sf z}%
_{4}\right)%
\end{array}%
\right)  \nonumber \\
&&\qquad \qquad \qquad \qquad \qquad \qquad \qquad \qquad \qquad \qquad
\qquad \qquad \times \left( 
\begin{array}{c}
Y_{K_{OD}}\left( {\sf z}_{3},{\sf z}_{4}\right) \\ 
Y_{K_{D}}\left( {\sf z}_{3},{\sf z}_{4}\right) \\ 
S_{K_{D}}\left( {\sf z}_{3},{\sf z}_{4}\right)%
\end{array}%
\right) d{\sf z}_{3}d{\sf z}_{4}  \label{kerneq}
\end{eqnarray}%
where 
\begin{eqnarray}
Y_{K_{OD}}\left( {\sf z}_{1,}{\sf z}_{2}\right) &=&\sum_{1\leq \iota \leq
\infty }\left( {\sf z}_{1,}{\sf z}_{2}|\breve{h}_{\iota }\right) \left( 
\breve{h}_{\iota }^{d}|\breve{Y}\right) =\int P_{K_{OD}}\left( {\sf z}_{1},%
{\sf z}_{2},{\sf z}_{3},{\sf z}_{4}\right) Y\left( {\sf z}_{3,}{\sf z}%
_{4}\right) d{\sf z}_{3}d{\sf z}_{4}  \nonumber \\
Y_{K_{D}}\left( {\sf z}_{1,}{\sf z}_{2}\right) &=&\sum_{1\leq \iota \leq
\infty }\left( {\sf z}_{1,}{\sf z}_{2}|\breve{k}_{\iota }\right) \left( 
\breve{k}_{\iota }^{d}|\breve{Y}\right) =\int P_{K_{D}}\left( {\sf z}_{1},%
{\sf z}_{2},{\sf z}_{3},{\sf z}_{4}\right) Y\left( {\sf z}_{3,}{\sf z}%
_{4}\right) d{\sf z}_{3}d{\sf z}_{4}  \nonumber \\
Y_{S_{D}}\left( {\sf z}_{1,}{\sf z}_{2}\right) &=&\sum_{1\leq \iota \leq
\infty }\left( {\sf z}_{1,}{\sf z}_{2}|\breve{s}_{\iota }\right) \left( 
\breve{s}_{\iota }^{d}|\breve{Y}\right) =\int P_{S_{D}}\left( {\sf z}_{1},%
{\sf z}_{2},{\sf z}_{3},{\sf z}_{4}\right) Y\left( {\sf z}_{3,}{\sf z}%
_{4}\right) d{\sf z}_{3}d{\sf z}_{4}
\end{eqnarray}

In the main body of this article we have considered maps, $\hat{T}$, that
map pure states to pure states and have the property $\hat{T}_{K_{OD}K_{D}}=%
\hat{T}_{K_{D}K_{OD}}=\hat{T}_{K_{OD}S_{D}}=0$. These transformations $\hat{T%
}$ are congruence transformations of ${\cal B}_{1}\left( {\cal H}^{N}\right) 
$ formed from multiplication operators (in the coordinate representation)
acting on ${\cal H}^{N}$ and act as {\it super multiplication operators }(in
the coordinate representation){\it \ }on ${\cal B}_{1}\left( {\cal H}%
^{N}\right) $ in the following fashion%
\begin{eqnarray}
\hat{T}\breve{X} &=&\breve{T}\breve{X}\breve{T}^{\dagger
}\Longleftrightarrow \text{ the integral kernel relationships}  \nonumber \\
&&\left( \left| {\sf z}_{1}\right\rangle \left\langle {\sf z}_{2}\right| |%
\hat{T}\breve{X}\right) =\left( \left| {\sf z}_{1}\right\rangle \left\langle 
{\sf z}_{2}\right| |\breve{T}\breve{X}\breve{T}^{\dagger }\right) =Tr\left\{
\left| {\sf z}_{2}\right\rangle \left\langle {\sf z}_{1}\right| \breve{T}%
\breve{X}\breve{T}^{\dagger }\right\}  \nonumber \\
\qquad &=&\left\langle {\sf z}_{1}|\breve{T}\breve{X}\breve{T}^{\dagger }%
{\sf z}_{2}\right\rangle =T\left( {\sf z}_{1}\right) X\left( {\sf z}_{1},%
{\sf z}_{2}\right) T^{\ast }\left( {\sf z}_{2}\right)  \nonumber \\
&=&\left( 1+\phi \left( {\sf z}_{1}\right) \right) ^{\frac{1}{2}}e^{i\omega
\left( {\sf z}_{1}\right) }X\left( {\sf z}_{1},{\sf z}_{2}\right) \left(
1+\phi \left( {\sf z}_{2}\right) \right) ^{\frac{1}{2}}e^{-i\omega \left( 
{\sf z}_{2}\right) }  \nonumber \\
&=&\left( 1+\phi \left( {\sf z}_{1}\right) \right) ^{\frac{1}{2}}\left(
1+\phi \left( {\sf z}_{2}\right) \right) ^{\frac{1}{2}}e^{i\left( \omega
\left( {\sf z}_{1}\right) -\omega \left( {\sf z}_{2}\right) \right) }X\left( 
{\sf z}_{1},{\sf z}_{2}\right)  \nonumber \\
&\equiv &{\rm T}\left( {\sf z}_{1},{\sf z}_{2}\right) X\left( {\sf z}_{1},%
{\sf z}_{2}\right)  \nonumber \\
&\Rightarrow &{\rm T}\left( {\sf z}_{1},{\sf z}_{2}\right) =\left( 1+\phi
\left( {\sf z}_{1}\right) \right) ^{\frac{1}{2}}\left( 1+\phi \left( {\sf z}%
_{2}\right) \right) ^{\frac{1}{2}}e^{i\left( \omega \left( {\sf z}%
_{1}\right) -\omega \left( {\sf z}_{2}\right) \right) }
\end{eqnarray}%
where, as in the main text, $T\left( {\sf z}_{1}\right) =\left( 1+\phi
\left( {\sf z}_{1}\right) \right) ^{\frac{1}{2}}e^{i\omega \left( {\sf z}%
_{1}\right) }$ $=e^{{\cal V}\left( {\sf z}_{1}\right) +i\omega \left( {\sf z}%
_{1}\right) }$ and hence the operators $\hat{T}$ can be expanded as 
\begin{eqnarray}
\hat{T} &=&\int {\rm T}\left( {\sf z}_{1},{\sf z}_{2}\right) \delta \left( 
{\sf z}_{1}-{\sf z}_{3}\right) \delta \left( {\sf z}_{2}-{\sf z}_{4}\right) |%
{\sf z}_{1},{\sf z}_{2}{\bf )}\left( {\sf z}_{3},{\sf z}_{4}\right| d{\sf z}%
_{1}d{\sf z}_{2}d{\sf z}_{3}d{\sf z}_{4}  \nonumber \\
&=&\int \left( 1+\phi \left( {\sf z}_{1}\right) \right) ^{\frac{1}{2}}\left(
1+\phi \left( {\sf z}_{2}\right) \right) ^{\frac{1}{2}}e^{i\left( \omega
\left( {\sf z}_{1}\right) -\omega \left( {\sf z}_{2}\right) \right) }|{\sf z}%
_{1},{\sf z}_{2}{\bf )}\left( {\sf z}_{1},{\sf z}_{2}\right| d{\sf z}_{1}d%
{\sf z}_{2}
\end{eqnarray}%
The operators $\hat{T}$ have the decomposition 
\begin{eqnarray}
\hat{T} &=&\hat{P}_{K_{OD}}\hat{T}\hat{P}_{K_{OD}}+\hat{P}_{K_{D}}\hat{T}%
\hat{P}_{K_{D}}+\hat{P}_{K_{D}}\hat{T}\hat{P}_{K_{S}}+\hat{P}_{S_{D}}\hat{T}%
\hat{P}_{S_{D}}  \nonumber \\
&\equiv &\hat{T}_{K_{OD}}+\hat{T}_{K_{D}}+\hat{T}_{K_{D}S_{D}}+\hat{T}%
_{S_{D}}
\end{eqnarray}%
one can see that $\hat{T}_{d}\equiv \hat{T}_{K_{D}}+\hat{T}_{K_{D}S_{D}}+%
\hat{T}_{K_{S}S_{D}}$ has the expansion 
\begin{equation}
\hat{T}_{d}=\int \left( 1+\phi \left( {\sf z}_{1}\right) \right) |{\sf z}%
_{1},{\sf z}_{1}{\bf )}\left( {\sf z}_{1},{\sf z}_{1}\right| d{\sf z}_{1}
\end{equation}%
We now consider $\hat{T}=\hat{I}+\hat{\Phi}$ and note that 
\begin{equation}
\hat{\Phi}_{d}=\int \phi \left( {\sf z}_{1}\right) |{\sf z}_{1},{\sf z}_{1}%
{\bf )}\left( {\sf z}_{1},{\sf z}_{1}\right| d{\sf z}_{1}
\end{equation}%
and the various projections and integral kernels of projections of $\hat{\Phi%
}_{d}$ are given by 
\begin{eqnarray}
\hat{\Phi}_{K_{D}} &=&\int \Phi _{K_{D}}\left( {\sf z}_{1},{\sf z}%
_{3}\right) |{\sf z}_{1},{\sf z}_{1}{\bf )}\left( {\sf z}_{3},{\sf z}%
_{3}\right| d{\sf z}_{1}d{\sf z}_{3}  \nonumber \\
\Phi _{K_{D}}\left( {\sf z}_{1},{\sf z}_{3}\right) &=&\int P_{K_{D}}\left( 
{\sf z}_{1},{\sf z}_{2}\right) \phi \left( {\sf z}_{2}\right)
P_{K_{D}}\left( {\sf z}_{2},{\sf z}_{3}\right) d{\sf z}_{2}  \nonumber \\
\hat{\Phi}_{K_{D}S_{D}} &=&\int \Phi _{K_{D}S_{D}}\left( {\sf z}_{1},{\sf z}%
_{3}\right) |{\sf z}_{1},{\sf z}_{1}{\bf )}\left( {\sf z}_{3},{\sf z}%
_{3}\right| d{\sf z}_{1}d{\sf z}_{3}  \nonumber \\
\Phi _{K_{D}S_{D}}\left( {\sf z}_{1},{\sf z}_{3}\right) &=&\int
P_{K_{D}}\left( {\sf z}_{1},{\sf z}_{2}\right) \phi \left( {\sf z}%
_{2}\right) P_{S_{D}}\left( {\sf z}_{2},{\sf z}_{3}\right) d{\sf z}_{2} 
\nonumber \\
\hat{\Phi}_{S_{D}} &=&\int \Phi _{S_{D}}\left( {\sf z}_{1},{\sf z}%
_{3}\right) |{\sf z}_{1},{\sf z}_{1}{\bf )}\left( {\sf z}_{3},{\sf z}%
_{3}\right| d{\sf z}_{1}d{\sf z}_{3}  \nonumber \\
\Phi _{S_{D}}\left( {\sf z}_{1},{\sf z}_{3}\right) &=&\int P_{S_{D}}\left( 
{\sf z}_{1},{\sf z}_{2}\right) \phi \left( {\sf z}_{2}\right)
P_{S_{D}}\left( {\sf z}_{2},{\sf z}_{3}\right) d{\sf z}_{2}  \label{kerproj}
\end{eqnarray}%
The super operator $\hat{T}$ can be associated with the following matrices
of integral kernels 
\begin{eqnarray}
&&\hat{T}{\frak \equiv }\left( 
\begin{array}{ccc}
\left( {\cal K}|\hat{T}_{K_{OD}}{\cal K}\right) & 0 & 0 \\ 
0 & \left( {\cal K}|\hat{T}_{K_{D}}{\cal K}\right) & \left( {\cal K}|\hat{T}%
_{K_{D}S_{D}}{\cal K}\right) \\ 
0 & 0 & \left( {\cal K}|\hat{T}_{S_{D}}{\cal K}\right)%
\end{array}%
\right) =  \nonumber \\
&&\left( 
\begin{array}{ccc}
\left( {\cal K}|\hat{P}_{K_{OD}}{\cal K}\right) & 0 & 0 \\ 
0 & \left( {\cal K}|\hat{P}_{K_{D}}{\cal K}\right) & 0 \\ 
0 & 0 & \left( {\cal K}|\hat{P}_{S_{D}}{\cal K}\right)%
\end{array}%
\right) +\left( 
\begin{array}{ccc}
\left( {\cal K}|\hat{\Phi}_{K_{OD}}{\cal K}\right) & 0 & 0 \\ 
0 & \left( {\cal K}|\hat{\Phi}_{K_{D}}{\cal K}\right) & \left( {\cal K}|\hat{%
\Phi}_{K_{D}S_{D}}{\cal K}\right) \\ 
0 & 0 & \left( {\cal K}|\hat{\Phi}_{S_{D}}{\cal K}\right)%
\end{array}%
\right)  \nonumber \\
&&
\end{eqnarray}%
The various integral kernels composing the transformed kernel evaluated in
the manner of Eqs.(\ref{transkern}) and (\ref{kerneq}) are 
\begin{equation}
Y^{\prime }\left( {\sf z}_{1},{\sf z}_{2}\right) =Y_{K_{OD}}^{\prime }\left( 
{\sf z}_{1},{\sf z}_{2}\right) +Y_{K_{D}}^{\prime }\left( {\sf z}_{1},{\sf z}%
_{2}\right) +Y_{S_{D}}^{\prime }\left( {\sf z}_{1},{\sf z}_{2}\right)
\end{equation}%
and are given by 
\begin{eqnarray}
Y_{K_{OD}}^{\prime }\left( {\sf z}_{1},{\sf z}_{2}\right) &=&\left( {\sf z}%
_{1},{\sf z}_{2}|\hat{T}\breve{Y}_{K_{OD}}\right) =\left( {\sf z}_{1},{\sf z}%
_{2}|\hat{T}_{K_{OD}}\breve{Y}_{K_{OD}}\right)  \nonumber \\
&=&\int \left\{ P_{K_{OD}}\left( {\sf z}_{1},{\sf z}_{2},{\sf z}_{3},{\sf z}%
_{4}\right) +\Phi _{K_{OD}}\left( {\sf z}_{1},{\sf z}_{2},{\sf z}_{3},{\sf z}%
_{4}\right) \right\} Y_{K_{OD}}\left( {\sf z}_{3},{\sf z}_{4}\right) d{\sf z}%
_{3}d{\sf z}_{4}  \nonumber \\
&=&\left( 1+\phi \left( {\sf z}_{1}\right) \right) ^{\frac{1}{2}}\left(
1+\phi \left( {\sf z}_{2}\right) \right) ^{\frac{1}{2}}e^{i\left( \omega
\left( {\sf z}_{1}\right) -\omega \left( {\sf z}_{2}\right) \right) }\left(
1-\delta \left( {\sf z}_{1}-{\sf z}_{2}\right) \right) Y_{K_{OD}}\left( {\sf %
z}_{1},{\sf z}_{2}\right)
\end{eqnarray}%
and 
\begin{eqnarray}
Y_{d}^{\prime }\left( {\sf z}_{1},{\sf z}_{1}\right) &=&\left( {\sf z}_{1},%
{\sf z}_{1}|\hat{T}\breve{Y}_{d}\right) =\left\langle {\sf z}_{1}|\breve{T}%
\breve{Y}_{d}\breve{T}^{\dagger }{\sf z}_{1}\right\rangle  \nonumber \\
&=&\int \left\{ P_{d}\left( {\sf z}_{1},{\sf z}_{1},{\sf z}_{3},{\sf z}%
_{3}\right) +\Phi _{K_{d}}\left( {\sf z}_{1},{\sf z}_{1},{\sf z}_{3},{\sf z}%
_{3}\right) \right\} Y_{d}\left( {\sf z}_{3},{\sf z}_{3}\right) d{\sf z}_{3}
\nonumber \\
&=&T\left( {\sf z}_{1}\right) Y_{d}\left( {\sf z}_{1},{\sf z}_{1}\right)
T^{\ast }\left( {\sf z}_{1}\right) =\left| T\left( {\sf z}_{1}\right)
\right| ^{2}Y_{d}\left( {\sf z}_{1}\right) =\left( 1+\phi \left( {\sf z}%
_{1}\right) \right) Y_{d}\left( {\sf z}_{1}\right)
\end{eqnarray}%
which splits into two parts $\left( {\sf z}_{1},{\sf z}_{1}|\hat{T}\breve{Y}%
_{d}\right) =\left( {\sf z}_{1},{\sf z}_{1}|\breve{Y}_{K_{D}}^{\prime
}\right) +\left( {\sf z}_{1},{\sf z}_{1}|\breve{Y}_{S_{D}}^{\prime }\right) $
given by 
\begin{eqnarray}
\left( {\sf z}_{1},{\sf z}_{1}|\breve{Y}_{K_{D}}^{\prime }\right) &=&\int
\left\{ \left( \delta \left( {\sf z}_{1}-{\sf z}_{2}\right) +\Phi
_{K_{D}}\left( {\sf z}_{1},{\sf z}_{2}\right) \right) Y_{K_{D}}\left( {\sf z}%
_{2}\right) +\Phi _{K_{D}S_{D}}\left( {\sf z}_{1},{\sf z}_{2}\right)
Y_{S_{D}}\left( {\sf z}_{2}\right) \right\} d{\sf z}_{2}  \nonumber \\
\left( {\sf z}_{1},{\sf z}_{1}|\breve{Y}_{S_{D}}^{\prime }\right) &=&\int
\left( \delta \left( {\sf z}_{1}-{\sf z}_{2}\right) +\Phi _{S_{D}}\left( 
{\sf z}_{1},{\sf z}_{2}\right) \right) Y_{S_{D}}\left( {\sf z}_{2}\right) d%
{\sf z}_{2}
\end{eqnarray}%
Collecting together the diagonal terms from the last two equations one
obtains 
\begin{eqnarray}
\left( {\sf z}_{1},{\sf z}_{1}|\hat{T}\breve{Y}_{d}\right) &=&\left( 1+\phi
\left( {\sf z}_{1}\right) \right) Y_{d}\left( {\sf z}_{1}\right)  \nonumber
\\
&=&\left\{ Y_{K_{D}}\left( {\sf z}_{1}\right) +Y_{S_{D}}\left( {\sf z}%
_{1}\right) \right\} +\int \left\{ \Phi _{K_{D}}\left( {\sf z}_{1},{\sf z}%
_{2}\right) Y_{K_{D}}\left( {\sf z}_{2}\right) \right.  \nonumber \\
&&+\left. \Phi _{K_{D}S_{D}}\left( {\sf z}_{1},{\sf z}_{2}\right)
Y_{S_{D}}\left( {\sf z}_{2}\right) +\Phi _{K_{D}S_{D}}\left( {\sf z}_{1},%
{\sf z}_{2}\right) Y_{S_{D}}\left( {\sf z}_{2}\right) \right\} d{\sf z}_{2}
\end{eqnarray}%
which gives the important relationship 
\begin{eqnarray}
\phi \left( {\sf z}_{1}\right) Y_{d}\left( {\sf z}_{1}\right) &=&\int
\left\{ \Phi _{K_{D}}\left( {\sf z}_{1},{\sf z}_{2}\right) Y_{K_{D}}\left( 
{\sf z}_{2}\right) \right.  \nonumber \\
&&+\left. \Phi _{K_{D}S_{D}}\left( {\sf z}_{1},{\sf z}_{2}\right)
Y_{S_{D}}\left( {\sf z}_{2}\right) +\Phi _{S_{D}}\left( {\sf z}_{1},{\sf z}%
_{2}\right) Y_{S_{D}}\left( {\sf z}_{2}\right) \right\} d{\sf z}_{2}
\end{eqnarray}%
The projected integral kernels of $\hat{\Phi}$ in the preceding equations
are given in Eq.$\left( \ref{kerproj}\right) $. To summarize, the action of
the $\hat{T}$ {\it super multiplication operator }(in the coordinate
representation) on an operator $\breve{Y}=\breve{Y}_{K_{OD}}+\breve{Y}_{d}=%
\breve{Y}_{K_{OD}}+\breve{Y}_{K_{D}}+\breve{Y}_{S_{D}}$, is given by%
\begin{eqnarray}
&&\left( 
\begin{array}{c}
\left( {\cal K}|\breve{Y}_{K_{OD}}^{\prime }\right) \\ 
\left( {\cal K}|\breve{Y}_{K_{D}}^{\prime }\right) \\ 
\left( {\cal K}|\breve{Y}_{S_{D}}^{\prime }\right)%
\end{array}%
\right) =\left( 
\begin{array}{ccc}
\left( {\cal K}|\hat{T}_{K_{OD}}{\cal K}\right) & 0 & 0 \\ 
0 & \left( {\cal K}|\hat{T}_{K_{D}}{\cal K}\right) & \left( {\cal K}|\hat{T}%
_{S_{D}K_{OD}}{\cal K}\right) \\ 
0 & 0 & \left( {\cal K}|\hat{T}_{S_{D}}{\cal K}\right)%
\end{array}%
\right) \left( 
\begin{array}{c}
\left( {\cal K}|\breve{Y}_{K_{OD}}\right) \\ 
\left( {\cal K}|\breve{Y}_{K_{D}}\right) \\ 
\left( {\cal K}|\breve{Y}_{S_{D}}\right)%
\end{array}%
\right) =  \nonumber \\
&&\left\{ \left( 
\begin{array}{ccc}
\left( {\cal K}|\hat{P}_{K_{OD}}{\cal K}\right) & 0 & 0 \\ 
0 & \left( {\cal K}|\hat{P}_{K_{D}}{\cal K}\right) & 0 \\ 
0 & 0 & \left( {\cal K}|\hat{P}_{S_{D}}{\cal K}\right)%
\end{array}%
\right) +\right.  \nonumber \\
&&\qquad \qquad \qquad \left. \left( 
\begin{array}{ccc}
\left( {\cal K}|\hat{\Phi}_{K_{OD}}{\cal K}\right) & 0 & 0 \\ 
0 & \left( {\cal K}|\hat{\Phi}_{K_{D}}{\cal K}\right) & \left( {\cal K}|\hat{%
\Phi}_{S_{D}K_{OD}}{\cal K}\right) \\ 
0 & 0 & \left( {\cal K}|\hat{\Phi}_{S_{D}}{\cal K}\right)%
\end{array}%
\right) \right\} \left( 
\begin{array}{c}
\left( {\cal K}|\breve{Y}_{K_{OD}}\right) \\ 
\left( {\cal K}|\breve{Y}_{K_{D}}\right) \\ 
\left( {\cal K}|\breve{Y}_{S_{D}}\right)%
\end{array}%
\right)
\end{eqnarray}%
and 
\begin{equation}
\left( {\cal K}|\breve{Y}^{\prime }\right) =\left( {\cal K}|\breve{Y}%
_{K_{OD}}^{\prime }\right) +\left( {\cal K}|\breve{Y}_{K_{D}}^{\prime
}\right) +\left( {\cal K}|\breve{Y}_{S_{D}}^{\prime }\right)
\end{equation}%
where 
\begin{eqnarray}
&&\left( {\sf z}_{1},{\sf z}_{1}|\breve{Y}_{K_{D}}^{\prime }\right) +\left( 
{\sf z}_{1},{\sf z}_{1}|\breve{Y}_{S_{D}}^{\prime }\right) =\left( 1+\phi
\left( {\sf z}_{1}\right) \right) Y_{d}\left( {\sf z}_{1}\right)  \nonumber
\\
&&\qquad \qquad \qquad \qquad =Y_{K_{D}}^{\prime }\left( {\sf z}_{1}\right)
+Y_{S_{D}}^{\prime }\left( {\sf z}_{1}\right)  \nonumber \\
&&\qquad \qquad \qquad \qquad +\Delta _{K_{D}}Y_{K_{D}}\left( {\sf z}%
_{1}\right) +\Delta _{K_{D}S_{D}}Y_{S_{D}}\left( {\sf z}_{1}\right) +\Delta
_{S_{D}}Y_{S_{D}}\left( {\sf z}_{1}\right)  \label{diagker}
\end{eqnarray}%
where the incremental change functions are 
\begin{eqnarray}
\Delta _{K_{D}}Y_{K_{D}}\left( {\sf z}_{1}\right) &=&\int \Phi
_{K_{D}}\left( {\sf z}_{1},{\sf z}_{2}\right) Y_{K_{D}}\left( {\sf z}%
_{2}\right) d{\sf z}_{2}  \nonumber \\
\Delta _{K_{D}S_{D}}Y_{S_{D}}\left( {\sf z}_{1}\right) &=&\int \Phi
_{K_{D}S_{D}}\left( {\sf z}_{1},{\sf z}_{2}\right) Y_{S_{D}}\left( {\sf z}%
_{2}\right) d{\sf z}_{2}  \nonumber \\
\Delta _{S_{D}}Y_{S_{D}}\left( {\sf z}_{1}\right) &=&\int \Phi
_{S_{D}}\left( {\sf z}_{1},{\sf z}_{2}\right) Y_{S_{D}}\left( {\sf z}%
_{2}\right) d{\sf z}_{2}  \label{incrfun}
\end{eqnarray}%
Note that in order for $\hat{T}$ $=I+\hat{\Phi}$ to be a vector bundle map :$%
{\frak B}_{vec}\rightarrow {\frak B}_{vec}$ the function $\phi $ defining
the integral kernel of $\hat{\Phi}$ must have the following properties 
\begin{equation}
\left( \breve{s}_{i}^{d}|\hat{\Phi}\breve{k}_{\kappa }\right) =\int s_{\iota
}^{d\ast }\left( {\sf z}\right) \phi \left( {\sf z}\right) k_{\kappa }\left( 
{\sf z}\right) d{\sf z}=0\quad \forall \iota ,\kappa
\end{equation}

\subsection{Intra and Inter Equivalence Class Transformations}

\subsubsection{Intra Equivalence Class Transformations}

In this section we utilize the decomposition of trace class operators 
\begin{equation}
\breve{Y}\left( f\right) =\breve{Y}_{K_{OD}}+\breve{Y}_{K_{D}}+\breve{\Gamma}%
\left( f\right)
\end{equation}%
adapted to the vector bundle structure of ${\frak B}_{vec}$, where $\Xi
\left( \breve{Y}\left( f\right) \right) =\Xi \left( \breve{\Gamma}\left(
f\right) \right) $ and $\Xi $ is an isomorphism from $S_{D}$ to $L_{1}\left( 
{\tt Z}\right) $. The equivalence class $\left[ f\right] _{N}$ is then given
by 
\begin{equation}
\left[ f\right] _{N}=\left\{ \breve{Y}_{K_{OD}}+\breve{Y}_{K_{D}}+\breve{%
\Gamma}\left( f\right) ;\breve{Y}_{K_{OD}}\in K_{OD},\breve{Y}_{K_{D}}\in
K_{D}\right\}
\end{equation}%
Transformations, $\hat{T}_{\ker },$ that map $\ker \Xi \rightarrow \ker \Xi $%
, where $\ker \Xi =K_{OD}\oplus K_{D}$, and have no effect on $S_{D\text{ }}$
are of the form 
\begin{eqnarray}
\hat{T}_{\ker } &=&\sum_{1\leq \kappa ,\iota \leq \infty }\left\{ {\rm T}%
_{\ker K_{OD}\kappa \iota }\left| \breve{h}_{\kappa }\right) \left( \breve{h}%
_{\iota }^{d}\right| +{\rm T}_{\ker K_{OD}K_{D}\kappa \iota }\left| \breve{h}%
_{\kappa }\right) \left( \breve{k}_{\iota }^{d}\right| \right.  \nonumber \\
&&\qquad \qquad \quad \left. +{\rm T}_{\ker K_{D}K_{OD}\kappa \iota }\left| 
\breve{k}_{\kappa }\right) \left( \breve{h}_{\iota }^{d}\right| +{\rm T}%
_{\ker K_{D}\kappa \iota }\left| \breve{k}_{\kappa }\right) \left( \breve{k}%
_{\iota }^{d}\right| +\left| \breve{s}_{\kappa }\right) \left( \breve{s}%
_{\iota }^{d}\right| \delta _{\iota \kappa }\right\} \\
&=&\hat{P}_{S_{D}}+\sum_{1\leq \kappa ,\iota \leq \infty }\left\{ {\rm T}%
_{\ker K_{OD}\kappa \iota }\left| \breve{h}_{\kappa }\right) \left( \breve{h}%
_{\iota }^{d}\right| +{\rm T}_{\ker K_{OD}K_{D}\kappa \iota }\left| \breve{h}%
_{\kappa }\right) \left( \breve{k}_{\iota }^{d}\right| \right.  \nonumber \\
&&\qquad \qquad \qquad \qquad \qquad \left. \,+{\rm T}_{\ker
K_{D}K_{OD}\kappa \iota }\left| \breve{k}_{\kappa }\right) \left( \breve{h}%
_{\iota }^{d}\right| +{\rm T}_{\ker K_{D}\kappa \iota }\left| \breve{k}%
_{\kappa }\right) \left( \breve{k}_{\iota }^{d}\right| \right\} \\
&=&\int {\rm T}_{\ker }\left( {\sf z}_{1},{\sf z}_{2},{\sf z}_{3},{\sf z}%
_{4}\right) |{\sf z}_{1},{\sf z}_{2}{\bf )}\left( {\sf z}_{3},{\sf z}%
_{4}\right| d{\sf z}_{1}d{\sf z}_{2}d{\sf z}_{3}d{\sf z}_{4}
\end{eqnarray}%
and map any given element of $\left[ f\right] _{N}$ into all others in $%
\left[ f\right] _{N}$.

Invertible transformations have the added property that they can be
expressed as exponentials 
\begin{equation}
\hat{T}_{\ker }=e^{\sum_{1\leq \kappa ,\iota \leq \infty }\left\{ {\rm %
\Upsilon }_{\ker K_{OD}\kappa \iota }\left| \breve{h}_{\kappa }\right)
\left( \breve{h}_{\iota }^{d}\right| +{\rm \Upsilon }_{\ker
K_{OD}K_{D}\kappa \iota }\left| \breve{h}_{\kappa }\right) \left( \breve{k}%
_{\iota }^{d}\right| +{\rm \Upsilon }_{\ker K_{D}K_{OD}\kappa \iota }\left| 
\breve{k}_{\kappa }\right) \left( \breve{h}_{\iota }^{d}\right| +{\rm %
\Upsilon }_{\ker K_{D}\kappa \iota }\left| \breve{k}_{\kappa }\right) \left( 
\breve{k}_{\iota }^{d}\right| \right\} }  \label{tker}
\end{equation}%
and form a representation of the group $GL\left( \ker \Xi \right) ,$ which
acts homogeneously\cite{homogen} in $\left[ f\right] _{N}$. More generally
we can form a representation of $GL\left( \ker \Xi \right) $ by invertible
super operators of the form 
\begin{equation}
\hat{T}_{\ker }^{\prime }=e^{\sum_{1\leq \iota ,\kappa \leq \infty }{\rm %
\Upsilon }_{\ker S_{D}}\left( \hat{T}_{\ker }\right) _{\kappa \iota }\left| 
\breve{s}_{\kappa }\right) \left( \breve{s}_{\iota }^{d}\right| }\hat{T}%
_{\ker }=\hat{t}\left( \hat{T}_{\ker }\right) \hat{T}_{\ker }
\end{equation}%
for coefficients $\left\{ {\rm \Upsilon }_{\ker S_{D}}\left( \hat{T}_{\ker
}\right) _{\kappa \iota }\right\} $ that are determined by the $\hat{T}%
_{\ker }$ given in Eq.$\left( \ref{tker}\right) $ and that satisfy the group
property%
\begin{equation}
\hat{t}\left( \hat{T}_{\ker 1}\right) \hat{t}\left( \hat{T}_{\ker 2}\right) =%
\hat{t}\left( \hat{T}_{\ker 1}\hat{T}_{\ker 2}\right)
\end{equation}%
but such transformations would not be intra fiber transformations. The
matrix coefficients $\left\{ {\rm \Upsilon }_{\ker K_{OD}\kappa \iota },{\rm %
\Upsilon }_{\ker K_{OD}K_{D}\kappa \iota },{\rm \Upsilon }_{\ker
K_{D}K_{OD}\kappa \iota },{\rm \Upsilon }_{\ker K_{D}\kappa \iota }\right\} $
parameterize the equivalence class in a redundant manner (i.e. two or more
sets of parameters are associated with a single element), however if we
start from a reference element of $\left[ f\right] _{N}$, which has the
decomposition, 
\begin{equation}
\breve{Y}_{ref}\left( f\right) =\breve{Y}_{refK_{OD}}+\breve{Y}_{refK_{D}}+%
\breve{\Gamma}\left( f\right)
\end{equation}%
we can generate all of $\left[ f\right] _{N}$ by acting on $\left| \breve{Y}%
_{ref}\left( f\right) \right) $ with the transformations 
\begin{equation}
\hat{T}_{\ker \breve{Y}_{ref}}=e^{\sum_{1\leq \kappa \leq \infty }\left\{
\Upsilon _{\ker K_{OD}\kappa \iota }\left| \breve{h}_{\kappa }\right) \left( 
\breve{Y}_{refK_{OD}}^{d}\right| +\Upsilon _{\ker K_{D}\kappa \iota }\left| 
\breve{k}_{\kappa }\right) \left( \breve{Y}_{refK_{D}}^{d}\right| \right\} }
\label{parker}
\end{equation}%
that form a coset as described in the summary section, where $\breve{h}_{1}=%
\breve{Y}_{refK_{OD}}$ and $\breve{k}_{1\iota }=\breve{Y}_{refK_{D}}$. This
gives a non redundant parameterization for the action of the group $GL\left(
\ker \Xi \right) $ on the element $\breve{Y}_{ref}\left( f\right) $, which
produces all the members of $\left[ f\right] _{N}$.

We now restrict attention to the subgroup, ${\frak G}_{fib}{\frak \subset }%
GL\left( \ker \Xi \right) $, of super multiplication operators, (in the
coordinate representation){\it \ }that can be represented by super operators
of the form $\hat{T}_{fib}=\hat{T}_{fibK_{OD}{\rm \ }}+\hat{T}_{fibK_{D}{\rm %
\ }}+\hat{P}_{S_{D}}$ where 
\begin{eqnarray}
\hat{T}_{fibK_{OD}{\rm \ }} &=&\hat{P}_{K_{OD}}+\hat{\Phi}_{K_{OD}} 
\nonumber \\
\hat{T}_{fibK_{D}{\rm \ }} &=&\hat{P}_{K_{D}}+\hat{\Phi}_{K_{D}}
\end{eqnarray}%
that map {\it pure states to pure states} and are generated by the functions 
$\phi _{K_{D}}$ and a fixed $\omega $ for $\hat{T}_{fibK_{OD}{\rm \ }}$,
that have the properties 
\begin{eqnarray}
\left( \breve{s}_{i}^{d}|\hat{\Phi}_{K_{D}}\breve{s}_{\kappa }\right)
&=&\int s_{\iota }^{d\ast }\left( {\sf z}\right) \phi _{K_{D}}\left( {\sf z}%
\right) s_{\kappa }\left( {\sf z}\right) d{\sf z}=0\quad \forall \iota
,\kappa  \nonumber \\
\left( \breve{k}_{i}^{d}|\hat{\Phi}_{K_{D}}\breve{s}_{\kappa }\right)
&=&\int k_{\iota }^{d\ast }\left( {\sf z}\right) \phi _{K_{D}}\left( {\sf z}%
\right) s_{\kappa }\left( {\sf z}\right) d{\sf z}=0\quad \forall \iota
,\kappa  \nonumber \\
\left( \breve{s}_{i}^{d}|\hat{\Phi}_{K_{D}}\breve{k}_{\kappa }\right)
&=&\int s_{\iota }^{d\ast }\left( {\sf z}\right) \phi _{K_{D}}\left( {\sf z}%
\right) k_{\kappa }\left( {\sf z}\right) d{\sf z}=0\quad \forall \iota
,\kappa  \nonumber \\
\left( \breve{k}_{i}^{d}|\hat{\Phi}_{K_{D}}\breve{k}_{\kappa }\right)
&=&\int k_{\iota }^{d\ast }\left( {\sf z}\right) \phi _{K_{D}}\left( {\sf z}%
\right) k_{\kappa }\left( {\sf z}\right) d{\sf z}\neq 0\quad \text{for some }%
\iota ,\kappa  \label{phicond}
\end{eqnarray}%
Functions with these properties can be expanded as 
\begin{eqnarray}
\phi _{K_{D}}\left( {\sf z}\right) &=&\sum_{1\leq \kappa ,\iota \leq \infty
}\phi _{K_{D}\kappa \iota }\left( {\sf z},{\sf z}|\breve{k}_{\kappa }\right)
\left( \breve{k}_{\iota }^{d}|{\sf z},{\sf z}\right)  \nonumber \\
&=&\int P_{K_{D}}\left( {\sf z},{\sf z}^{\prime }\right) \phi _{K_{D}}\left( 
{\sf z}^{\prime }\right) P_{K_{D}}\left( {\sf z}^{\prime },{\sf z}\right) d%
{\sf z}^{\prime }  \nonumber \\
\phi _{K_{D}\kappa \iota } &=&\int k_{\kappa }^{d\ast }\left( {\sf z}\right)
\phi _{K_{D}}\left( {\sf z}\right) k_{\iota }\left( {\sf z}\right) d{\sf z}
\label{PhiKD}
\end{eqnarray}%
However note that it is {\bf not }sufficient for functions to have
expansions of the form Eq.$\left( \ref{PhiKD}\right) $ in order for them to
satisfy the conditions expressed in Eq.$\left( \ref{phicond}\right) $, it is
only necessary. This can be seen, for instance, by considering the following
expansions with arbitrary $\phi _{K_{D}\kappa \iota }$ 
\begin{eqnarray}
&&\int s_{r}^{d\ast }\left( {\sf z}\right) \sum_{1\leq \kappa ,\iota \leq
\infty }\phi _{K_{D}\kappa \iota }\left( {\sf z},{\sf z}|\breve{k}_{\kappa
}\right) \left( \breve{k}_{\iota }^{d}|{\sf z},{\sf z}\right) s_{s}\left( 
{\sf z}\right) d{\sf z}  \nonumber \\
&=&\int \sum_{1\leq \kappa ,\iota \leq \infty }\phi _{K_{D}\kappa \iota
}\left( \breve{s}_{r}^{d}|{\sf z},{\sf z}\right) \left( {\sf z},{\sf z}|%
\breve{k}_{\kappa }\right) \left( \breve{k}_{\iota }^{d}|{\sf z},{\sf z}%
\right) \left( {\sf z},{\sf z}|\breve{s}_{s}\right) d{\sf z}
\end{eqnarray}%
which {\it might or might not} equal zero. The super operator $\hat{\Phi}%
_{K_{D}}$ is defined in terms of functions satisfying Eq.$\left( \ref%
{phicond}\right) $by the integral kernel 
\begin{equation}
\Phi _{K_{D}}\left( {\sf z}_{1,}{\sf z}_{2}\right) =\phi _{K_{D}}\left( {\sf %
z}_{1}\right) \delta \left( {\sf z}_{1}-{\sf z}_{2}\right)
\end{equation}%
The diagonal part of the transformed integral kernel $\left( {\sf z}_{1},%
{\sf z}_{1}|\breve{Y}_{d}^{\prime }\right) =\left( {\sf z}_{1},{\sf z}_{1}|%
\hat{T}_{fib}\breve{Y}_{d}\right) $ can be developed as 
\begin{eqnarray}
&&\left( {\sf z}_{1},{\sf z}_{1}|\breve{Y}_{K_{D}}^{\prime }\right) +\left( 
{\sf z}_{1},{\sf z}_{1}|\breve{\Gamma}\left( f\right) \right) =  \nonumber \\
&&\qquad \qquad \left( 1+\phi _{K_{D}}\left( {\sf z}_{1}\right) \right)
Y_{d}\left( {\sf z}_{1}\right) =  \nonumber \\
&&Y_{K_{D}}\left( {\sf z}_{1}\right) +\Gamma \left( f\right) \left( {\sf z}%
_{1}\right) +\Delta _{K_{D}}Y_{K_{D}}\left( {\sf z}_{1}\right) +\Delta
_{K_{D}S_{D}}\Gamma \left( f\right) \left( {\sf z}_{1}\right) +\Delta
_{S_{D}}\Gamma \left( f\right) \left( {\sf z}_{1}\right)  \nonumber \\
&&\qquad \qquad \qquad \qquad \qquad \qquad \quad =Y_{K_{D}}\left( {\sf z}%
_{1}\right) +\Gamma \left( f\right) \left( {\sf z}_{1}\right) +\Delta
_{K_{D}}Y_{K_{D}}\left( {\sf z}_{1}\right)  \label{diagexp}
\end{eqnarray}%
the last line in the above equation results from 
\begin{eqnarray}
\Phi _{K_{D}S_{D}}\left( {\sf z}_{1},{\sf z}_{2}\right) &=&\int
P_{K_{D}}\left( {\sf z}_{1},{\sf z}_{1}^{\prime }\right) \phi _{K_{D}}\left( 
{\sf z}_{1}^{\prime }\right) P_{S_{D}}\left( {\sf z}_{1}^{\prime },{\sf z}%
_{2}\right) d{\sf z}_{1}^{\prime }=0  \label{KDSD} \\
\Phi _{S_{D}}\left( {\sf z}_{1},{\sf z}_{2}\right) &=&\int P_{S_{D}}\left( 
{\sf z}_{1},{\sf z}_{1}^{\prime }\right) \phi _{K_{D}}\left( {\sf z}%
_{1}^{\prime }\right) P_{S_{D}}\left( {\sf z}_{1}^{\prime },{\sf z}%
_{2}\right) d{\sf z}_{1}^{\prime }=0  \label{PhiSD}
\end{eqnarray}%
and the definitions of the incremental functions Eq.$\left( \ref{incrfun}%
\right) $. Thus Eq.$\left( \ref{diagexp}\right) $ provides the identity 
\begin{equation}
\phi _{K_{D}}\left( {\sf z}_{1}\right) Y_{d}\left( {\sf z}_{1}\right)
=\Delta _{K_{D}}Y_{K_{D}}\left( {\sf z}_{1}\right) =\phi _{K_{D}}\left( {\sf %
z}_{1}\right) Y_{K_{D}}\left( {\sf z}_{1}\right)
\end{equation}%
as a result of Eqs.$\left( \ref{KDSD}\right) $, $\left( \ref{PhiSD}\right) $
and 
\begin{eqnarray}
\Delta _{K_{D}}Y_{K_{D}}\left( {\sf z}_{1}\right) &=&\int \Phi
_{K_{D}}\left( {\sf z}_{1},{\sf z}_{2}\right) Y_{K_{D}}\left( {\sf z}%
_{2}\right) d{\sf z}_{2}  \nonumber \\
&=&\int \phi _{K_{D}}\left( {\sf z}_{1}\right) \delta \left( {\sf z}_{1}-%
{\sf z}_{2}\right) Y_{K_{D}}\left( {\sf z}_{2}\right) d{\sf z}_{2}=\phi
_{K_{D}}\left( {\sf z}_{1}\right) Y_{K_{D}}\left( {\sf z}_{1}\right)
\end{eqnarray}%
The integral kernel of the super operator $\hat{T}_{fibK_{D}{\rm \ }}$ thus
has the expansions 
\begin{eqnarray}
\hat{T}_{fibK_{D}{\rm \ }} &=&\hat{P}_{K_{D}}+\hat{\Phi}_{K_{D}} \\
&=&\int \left\{ P_{K_{D}}\left( {\sf z}_{1},{\sf z}_{2}\right) +\phi
_{K_{D}}\left( {\sf z}_{1}\right) \delta \left( {\sf z}_{1}-{\sf z}%
_{2}\right) \right\} \left| {\sf z}_{1},{\sf z}_{2}\right) \left( {\sf z}%
_{1},{\sf z}_{2}\right| d{\sf z}_{1}d{\sf z}_{2}
\end{eqnarray}%
The integral kernel of the super operator $\hat{T}_{fibK_{OD}{\rm \ }}=\hat{P%
}_{K_{OD}}+\hat{\Phi}_{K_{OD}}$ which is determined up to phase by $\phi
_{K_{D}}\left( {\sf z}_{1}\right) $ has the expansion 
\begin{equation}
\Phi _{K_{OD}}\left( {\sf z}_{1},{\sf z}_{2},{\sf z}_{3,}{\sf z}_{4}\right)
=\sum_{1\leq \kappa ,\iota \leq \infty }\Phi _{K_{OD}\kappa \iota }\left( 
{\sf z}_{1},{\sf z}_{2}|\breve{h}_{\kappa }\right) \left( \breve{h}_{\iota
}^{d}|{\sf z}_{3},{\sf z}_{4}\right)
\end{equation}%
where 
\begin{equation}
\Phi _{K_{OD}\kappa \iota }=\int k_{\kappa }^{d\ast }\left( {\sf z}_{1},{\sf %
z}_{2}\right) \Phi _{K_{OD}}\left( {\sf z}_{1,}{\sf z}_{2},{\sf z}_{3,}{\sf z%
}_{4}\right) k_{\iota }\left( {\sf z}_{3},{\sf z}_{4}\right) d{\sf z}_{1}d%
{\sf z}_{2}d{\sf z}_{3}d{\sf z}_{4}
\end{equation}%
and has the projection properties 
\begin{eqnarray}
&&\Phi _{K_{OD}}\left( {\sf z}_{1,}{\sf z}_{2},{\sf z}_{3,}{\sf z}_{4}\right)
\nonumber \\
&=&\int P_{K_{OD}}\left( {\sf z}_{1},{\sf z}_{2},{\sf z}_{1}^{\prime },{\sf z%
}_{2}^{\prime }\right) \Phi _{K_{OD}}\left( {\sf z}_{1}^{\prime },{\sf z}%
_{2}^{\prime },{\sf z}_{3}^{\prime },{\sf z}_{4}^{\prime }\right)
P_{K_{OD}}\left( {\sf z}_{3}^{\prime },{\sf z}_{4}^{\prime },{\sf z}_{3},%
{\sf z}_{4}\right) d{\sf z}_{1}^{\prime }d{\sf z}_{2}^{\prime }d{\sf z}%
_{3}^{\prime }d{\sf z}_{4}^{\prime }  \nonumber \\
&=&\int P_{K_{OD}}\left( {\sf z}_{1},{\sf z}_{2},{\sf z}_{1}^{\prime },{\sf z%
}_{2}^{\prime }\right) \Phi _{K_{OD}}\left( {\sf z}_{1}^{\prime },{\sf z}%
_{2}^{\prime },{\sf z}_{3}^{\prime },{\sf z}_{4}^{\prime }\right) d{\sf z}%
_{1}^{\prime }d{\sf z}_{2}^{\prime }  \nonumber \\
&=&\int \Phi _{K_{OD}}\left( {\sf z}_{1},{\sf z}_{2},{\sf z}_{3}^{\prime },%
{\sf z}_{4}^{\prime }\right) P_{K_{OD}}\left( {\sf z}_{3}^{\prime },{\sf z}%
_{4}^{\prime },{\sf z}_{3},{\sf z}_{4}\right) d{\sf z}_{3}^{\prime }d{\sf z}%
_{4}^{\prime }
\end{eqnarray}%
The action of $\hat{T}_{fibK_{OD}}$ on $\hat{Y}_{K_{OD}}$ is given by 
\begin{eqnarray}
&&Y_{K_{OD}}^{\prime }\left( {\sf z}_{1,}{\sf z}_{2}\right) =\int
T_{fibK_{OD}}\left( {\sf z}_{1,}{\sf z}_{2},{\sf z}_{3},{\sf z}_{4}\right)
Y_{K_{OD}}\left( {\sf z}_{3,}{\sf z}_{4}\right) d{\sf z}_{3}d{\sf z}_{4} 
\nonumber \\
&=&\int \left\{ P_{K_{OD}}\left( {\sf z}_{1,}{\sf z}_{2},{\sf z}_{3},{\sf z}%
_{4}\right) +\Phi _{K_{OD}}\left( {\sf z}_{1,}{\sf z}_{2},{\sf z}_{3},{\sf z}%
_{4}\right) \right\} Y_{K_{OD}}\left( {\sf z}_{3},{\sf z}_{4}\right) d{\sf z}%
_{3}d{\sf z}_{4}  \nonumber \\
&=&\int \left( 1+\phi _{K_{D}}\left( {\sf z}_{1}\right) \right) ^{\frac{1}{2}%
}\left( 1+\phi _{K_{D}}\left( {\sf z}_{3}\right) \right) ^{\frac{1}{2}%
}e^{i\left( \omega \left( {\sf z}_{1}\right) -\omega \left( {\sf z}%
_{3}\right) \right) }  \nonumber \\
&&\times \left( 1-\delta \left( {\sf z}_{1}-{\sf z}_{2}\right) \right)
\left( 1-\delta \left( {\sf z}_{3}-{\sf z}_{4}\right) \right) \delta \left( 
{\sf z}_{1}-{\sf z}_{3}\right) \delta \left( {\sf z}_{3}-{\sf z}_{4}\right)
Y_{K_{OD}}\left( {\sf z}_{3},{\sf z}_{4}\right) d{\sf z}_{3}d{\sf z}_{4} 
\nonumber \\
&=&\left( 1+\phi _{K_{D}}\left( {\sf z}_{1}\right) \right) ^{\frac{1}{2}%
}\left( 1+\phi _{K_{D}}\left( {\sf z}_{2}\right) \right) ^{\frac{1}{2}%
}\left( 1-\delta \left( {\sf z}_{1}-{\sf z}_{2}\right) \right) e^{i\left(
\omega \left( {\sf z}_{1}\right) -\omega \left( {\sf z}_{2}\right) \right)
}Y_{K_{OD}}\left( {\sf z}_{1},{\sf z}_{2}\right)  \label{Yprime}
\end{eqnarray}%
where the phase $\omega $ can be {\it arbitrarily} chosen and represents the
factor from the subgroup ${\cal G}_{\Xi }$. The factor multiplying $%
Y_{K_{OD}}\left( {\sf z}_{1,}{\sf z}_{2}\right) $ in the last line of Eq.$%
\left( \ref{Yprime}\right) $ can be expanded further to give%
\begin{equation}
Eq.\left( \ref{Yprime}\right) =\left( 1+\Lambda _{_{K_{D}}}\left( {\sf z}%
_{1,}{\sf z}_{2}\right) \right) e^{i\left( \omega \left( {\sf z}_{1}\right)
-\omega \left( {\sf z}_{2}\right) \right) }\left( 1-\delta \left( {\sf z}%
_{1}-{\sf z}_{2}\right) \right)  \label{xpan}
\end{equation}%
where $\Lambda _{_{K_{D}}}\left( {\sf z}_{1,}{\sf z}_{2}\right) $ is a power
series in $\left( \phi _{K_{D}}\left( {\sf z}_{1}\right) +\phi
_{K_{D}}\left( {\sf z}_{2}\right) +\phi _{K_{D}}\left( {\sf z}_{1}\right)
\phi _{K_{D}}\left( {\sf z}_{2}\right) \right) $. Hence one can make the
identification 
\begin{equation}
\Phi _{K_{OD}}\left( {\sf z}_{1,}{\sf z}_{2},{\sf z}_{3},{\sf z}_{4}\right)
=\Lambda _{_{K_{D}}}\left( {\sf z}_{1,}{\sf z}_{2}\right) \delta \left( {\sf %
z}_{1}-{\sf z}_{3}\right) \delta \left( {\sf z}_{2}-{\sf z}_{4}\right)
\left( 1-\delta \left( {\sf z}_{1}-{\sf z}_{2}\right) \right)
\end{equation}%
To recap the complete super operator $\hat{T}_{fib}$ is given by 
\begin{equation}
\hat{T}_{fib}=\hat{T}_{fibK_{OD}{\rm \ }}+\hat{T}_{fibK_{D}{\rm \ }}+\hat{P}%
_{S_{D}}
\end{equation}%
where the various components are defined by the integral kernels in the
preceding equations. One should note that in contrast to $GL\left( \ker \Xi
\right) $ the subgroup ${\frak G}_{fib}$ does not act homogeneously on $%
\left[ f\right] _{N}$.

A non redundant parameterization for the action of the subgroup ${\frak G}%
_{fib}$ on a reference operator $\left| \breve{k}_{ref}\right) $ is obtained
by replacing $\phi _{K_{D}}\left( {\sf z}_{1}\right) $ in Eq.$\left( \ref%
{PhiKD}\right) $ by functions that have the additional properties 
\begin{equation}
\left( \breve{k}_{i}^{d}|\hat{\Phi}_{K_{D}}\breve{k}_{\kappa }\right) =\int
k_{\iota }^{d\ast }\left( {\sf z}\right) \phi _{K_{D}}\left( {\sf z}\right)
k_{\kappa }\left( {\sf z}\right) d{\sf z}=0;\quad 2\leq \kappa \leq \infty
,1\leq \iota \leq \infty  \label{phinr}
\end{equation}%
where $\left| \breve{k}_{1}\right) \equiv \left| \breve{k}_{ref}\right) =%
\hat{P}_{K_{D}}\left| \breve{Y}_{ref}\left( f\right) \right) $ to giving the
transformation $\hat{T}_{fib}\left( \breve{k}_{ref}\right) $. The functions
in Eq.$\left( \ref{phinr}\right) $ have the expansion 
\begin{eqnarray}
\phi _{refK_{D}}\left( {\sf z}_{1}\right) &=&\sum_{2\leq \kappa \leq \infty
}\phi _{refK_{D\kappa }}\left( {\sf z}_{1},{\sf z}_{1}|\breve{k}_{\kappa
}\right) \left( \breve{k}_{ref}^{d}|{\sf z}_{1},{\sf z}_{1}\right)  \nonumber
\\
\phi _{refK_{D\kappa }} &=&\int k_{\kappa }^{d\ast }\left( {\sf z}%
_{1}\right) \phi _{K_{D}}\left( {\sf z}_{1}\right) k_{ref}\left( {\sf z}%
_{1}\right) d{\sf z}_{1};\quad 2\leq \kappa \leq \infty  \label{refkd}
\end{eqnarray}%
Unlike the non redundant parameterization of $GL\left( \ker \Xi \right) $,
which produced all of $\left[ f\right] _{N}$ from the reference operator $%
\breve{Y}_{ref}\left( f\right) ,$ ${\frak G}_{fib}$ {\it does not. }However
if $\breve{Y}_{ref}\left( \gamma \right) $ is a pure state this
parameterization of {\it \ }${\frak G}_{fib}$ {\it does }produce all pure
states belonging to $\left[ \gamma \right] _{N}$ when acting on $\left| 
\breve{Y}_{ref}\left( \gamma \right) \right) \forall \gamma \in L_{1}\left( 
{\tt Z}\right) _{+}$.

\subsubsection{Inter Equivalence Class Transformations}

We now wish to consider transformations that map between equivalence
classes, and in particular to find a non redundant parameterization of $%
S_{D} $ by such transformations acting on a reference element. The most
obvious source of such transformations is the subgroup $GL\left(
S_{D}\right) $ represented by super operators of the form 
\begin{eqnarray}
\hat{T}_{S}=e^{\sum_{1\leq \kappa ,\iota \leq \infty }\Upsilon _{S_{D}\kappa
\iota }\left| \breve{s}_{\kappa }\right) \left( \breve{s}_{\iota
}^{d}\right| } &=&\hat{P}_{K_{OD}}+\hat{P}_{K_{D}}+\hat{t}_{S}  \nonumber \\
\hat{t}_{S} &=&\hat{P}_{S_{D}}\hat{T}_{S}\hat{P}_{S_{D}}
\end{eqnarray}%
More generally, however, we can form a reducible representation of $GL\left(
S_{D}\right) $ by invertible super operators of the form 
\begin{eqnarray}
\hat{T}_{S} &=&\hat{T}_{K_{OD}}\left( \hat{t}_{S}\right) +\hat{P}_{K_{D}}+%
\hat{t}_{S}  \nonumber \\
\hat{T}_{K_{OD}}\left( \hat{t}_{S}\right) &=&\hat{P}_{K_{OD}}\hat{T}_{S}\hat{%
P}_{K_{OD}}
\end{eqnarray}%
which are both inter fiber and intra fiber, where $\hat{T}_{K_{OD}}$ is a
representation of $GL\left( S_{D}\right) $. Considering these more general
representations of $GL\left( S_{D}\right) $ allows us to construct non
redundant parameterizations of $S_{D}$ using the a coset of the subgroup $%
{\frak G}_{den}\subset $ $GL\left( S_{D}\right) $ of super multiplication
operators (in the coordinate representation),{\it \ }%
\begin{equation}
\hat{T}_{den}\left( \phi _{S_{D}}\right) =\hat{T}_{K_{OD}{\rm \ }}\left(
\phi _{S_{D}},\omega \right) +\hat{P}_{K_{D}}+\hat{T}_{S_{D}{\rm \ }}\left(
\phi _{S_{D}}\right)  \label{tden}
\end{equation}%
generated by a $\phi _{S_{D}}$, where $\omega $ can be chosen to be zero,
that have the properties 
\begin{eqnarray}
\left( \breve{s}_{i}^{d}|\hat{\Phi}_{S_{D}}\breve{s}_{ref}\right) &=&\int
s_{\iota }^{d\ast }\left( {\sf z}\right) \phi _{S_{D}}\left( {\sf z}\right)
s_{ref}\left( {\sf z}\right) d{\sf z}\neq 0;\quad \text{for some }\iota
;1\leq \iota \leq \infty  \nonumber \\
\left( \breve{s}_{i}^{d}|\hat{\Phi}_{S_{D}}\breve{s}_{\kappa }\right)
&=&\int s_{\iota }^{d\ast }\left( {\sf z}\right) \phi _{S_{D}}\left( {\sf z}%
\right) s_{\kappa }\left( {\sf z}\right) d{\sf z}=0;\quad 1\leq \iota \leq
\infty ,2\leq \kappa \leq \infty  \nonumber \\
\left( \breve{k}_{i}^{d}|\hat{\Phi}_{S_{D}}\breve{s}_{\kappa }\right)
&=&\int k_{\iota }^{d\ast }\left( {\sf z}\right) \phi _{S_{D}}\left( {\sf z}%
\right) s_{\kappa }\left( {\sf z}\right) d{\sf z}=0;\quad 1\leq \iota
,\kappa \leq \infty  \nonumber \\
\left( \breve{s}_{i}^{d}|\hat{\Phi}_{S_{D}}\breve{k}_{\kappa }\right)
&=&\int s_{\iota }^{d\ast }\left( {\sf z}\right) \phi _{S_{D}}\left( {\sf z}%
\right) k_{\kappa }\left( {\sf z}\right) d{\sf z}=0;\quad 1\leq \iota
,\kappa \leq \infty  \nonumber \\
\left( \breve{k}_{i}^{d}|\hat{\Phi}_{S_{D}}\breve{k}_{\kappa }\right)
&=&\int k_{\iota }^{d\ast }\left( {\sf z}\right) \phi _{S_{D}}\left( {\sf z}%
\right) k_{\kappa }\left( {\sf z}\right) d{\sf z}=0;\quad 1\leq \iota
,\kappa \leq \infty
\end{eqnarray}%
{\it together} with a fixed choice of phase $\omega $ for $\hat{T}_{K_{OD}%
{\rm \ }}\left( \phi _{S_{D}}\right) $, where $\breve{s}_{1}=\breve{\Gamma}%
\left( f_{ref}\right) =\hat{P}_{S_{D}}\breve{Y}_{ref}$ \ and $\Xi \left( 
\breve{Y}_{ref}\right) =f_{ref}$. This coset is described in the summary
section. The function $\phi _{S_{D}}$ has the expansion and properties 
\begin{eqnarray}
\phi _{S_{D}}\left( {\sf z}\right) &=&\sum_{1\leq \kappa \leq \infty }\phi
_{S_{D}\kappa ref}\left( {\sf z},{\sf z}|\breve{s}_{\kappa }\right) \left( 
\breve{\Gamma}\left( f_{ref}\right) ^{d}|{\sf z},{\sf z}\right)  \nonumber \\
&=&\int P_{S_{D}}\left( {\sf z},{\sf z}^{\prime }\right) \phi _{S_{D}}\left( 
{\sf z}^{\prime }\right) P_{S_{D}}\left( {\sf z}^{\prime },{\sf z}\right) d%
{\sf z}^{\prime }  \nonumber \\
\phi _{S_{D}\kappa ref} &=&\int s_{\kappa }^{d\ast }\left( {\sf z}\right)
\phi _{S_{D}}\left( {\sf z}\right) \breve{\Gamma}\left( f_{ref}\right)
\left( {\sf z}\right) d{\sf z}
\end{eqnarray}%
The super operators $\hat{T}_{den}\left( \phi _{S_{D}}\right) $ generate
reference elements for all equivalence classes in a unique fashion and form
a non redundant parameterization for the action of the subgroup ${\frak G}%
_{den}$ {\it for a fixed choice of phase function }$\omega $ i.e. 
\begin{equation}
\hat{T}_{den}\left( \phi _{S_{D}}\right) :\left[ f_{ref}\right]
_{N}\rightarrow \left[ f^{\prime }\right] _{N};\quad f\neq f^{\prime }
\end{equation}%
The integral kernels of the super operators appearing in Eq.$\left( \ref%
{tden}\right) $, with suitable replacements of $\phi _{K_{D}}$ by $\phi
_{S_{D}}$, are given by%
\begin{eqnarray}
T_{K_{OD}{\rm \ }}\left( \phi _{S_{D}}\right) \left( {\sf z}_{1,}{\sf z}_{2},%
{\sf z}_{3},{\sf z}_{4}\right) &=&P_{K_{OD}}\left( {\sf z}_{1,}{\sf z}_{2},%
{\sf z}_{3},{\sf z}_{4}\right) +\Phi _{K_{OD}}\left( {\sf z}_{1,}{\sf z}_{2},%
{\sf z}_{3},{\sf z}_{4}\right)  \nonumber \\
\Phi _{K_{OD}}\left( {\sf z}_{1,}{\sf z}_{2},{\sf z}_{3},{\sf z}_{4}\right)
&=&\Lambda _{_{K_{S}}}\left( {\sf z}_{1},{\sf z}_{2}\right) \delta \left( 
{\sf z}_{1}-{\sf z}_{3}\right) \delta \left( {\sf z}_{2}-{\sf z}_{4}\right)
\left( 1-\delta \left( {\sf z}_{1}-{\sf z}_{2}\right) \right)  \label{tden2}
\end{eqnarray}%
(notice the subscript $K_{S\text{ }}$ in the definition of $\Lambda
_{_{K_{S}}}$) and 
\begin{equation}
T_{S_{D}{\rm \ }}\left( \phi _{K_{S}}\right) \left( {\sf z}_{1,}{\sf z}%
_{2}\right) =P_{K_{S}}\left( {\sf z}_{1,}{\sf z}_{2}\right) +\phi
_{K_{S}}\left( {\sf z}_{1}\right) \delta \left( {\sf z}_{1}-{\sf z}%
_{2}\right)  \label{tden3}
\end{equation}%
If $\left| \breve{Y}_{ref}\right) $ is a pure state then $\hat{T}%
_{den}\left( \phi _{S_{D}}\right) \left| \breve{Y}_{ref}\right) $ is also a
pure state different from $\left| \breve{Y}_{ref}\right) $. The set $\left\{ 
\hat{T}_{den}\left( \phi _{S_{D}}\right) \breve{Y}_{ref};\quad \forall \phi
_{S_{D}},\text{ fixed }\omega \right\} $ does not produce all pure states 
{\it but it does contain a unique pure state representative for each
equivalence class. }{\bf In the vector bundle classification of }${\cal B}%
_{1}\left( {\cal H}^{N}\right) $ {\bf these different} {\bf paths of pure
state represenatives provide different representations of the base of the
vector bundle }${\frak B}_{vect}$.

\section{Summary}

We now summarize the salient features of this document and relate them to
ideas introduced in the main body of paper I. The group of transformation, $%
{\frak G}_{\Xi }$, that preserves the vector bundle structure, ${\frak B}%
_{vect}$, of ${\cal B}_{1}\left( {\cal H}^{N}\right) $ is composed of maps
that have the decomposition 
\begin{equation}
{\frak \hat{T}}{\frak =}{\frak \hat{T}}_{K_{OD}}+{\frak \hat{T}}_{K_{D}}+%
{\frak \hat{T}}_{S_{D}}+{\frak \hat{T}}_{K_{D}S_{D}}+{\frak \hat{T}}%
_{K_{OD}S_{D}}+{\frak \hat{T}}_{K_{OD}K_{D}}+{\frak \hat{T}}_{K_{D}K_{OD}}
\end{equation}%
and refer to the direct decomposition of ${\cal B}_{1}\left( {\cal H}%
^{N}\right) $%
\begin{equation}
{\cal B}_{1}\left( {\cal H}^{N}\right) =\ker \Xi \oplus S_{D}=K_{OD}\oplus
K_{D}\oplus S_{D}
\end{equation}%
The subgroup $GL\left( \ker \Xi \right) $ maps each equivalence class to
itself and are thus intra fiber transformations, while the subgroup $%
GL\left( S_{D}\right) $, {\it in general},{\it \ }maps between equivalence
classes, but contains the subgroup, ${\frak G}_{den}$, of inter fiber
transformations. These subgroups can be used to give non redundant
parameterizations of $\ker \Xi $ and $S_{D}\equiv $ $\ker \Xi ^{c}$ by
choosing a reference state $\breve{Y}_{ref}\in {\cal B}_{1}\left( {\cal H}%
^{N}\right) $, where $\Xi \left( \breve{Y}_{ref}\right) =f_{ref}$, which can
be decomposed as 
\begin{eqnarray}
Y_{ref} &=&Y_{ref\,K_{OD}}+Y_{ref\,K_{D}}+\breve{\Gamma}\left( f_{ref}\right)
\\
\left| \breve{\Gamma}\left( f_{ref}\right) \right) &=&\hat{P}_{S_{D}}\left| 
\breve{Y}_{ref}\right)  \label{ref}
\end{eqnarray}%
and then forming left cosets, ${\frak Co}_{\ker \Xi }$ and ${\frak Co}%
_{S_{D}}$ of $GL\left( \ker \Xi \right) $ and $GL\left( S_{D}\right) $ with
respect to the isotropy subgroups of $\left| \breve{Y}_{ref\,K_{D}}\right) $
and $\left| \breve{\Gamma}\left( f_{ref}\right) \right) $ respectively in
the following manner. 
\begin{eqnarray}
{\frak G}_{Y_{ref\,d}} &=&\left\{ {\frak \hat{T};\hat{T}}\left| \breve{Y}%
_{ref\,K_{D}}\right) =\left| \breve{Y}_{ref\,K_{D}}\right) \right\} 
\nonumber \\
{\frak Co}_{\ker \Xi } &=&GL\left( \ker \Xi \right) /{\frak G}_{\breve{Y}%
_{ref\,K_{D}}} \\
{\frak G}_{\Gamma \left( f_{ref}\right) } &=&\left\{ {\frak \hat{T};\hat{T}}%
\left| \breve{\Gamma}\left( f_{ref}\right) \right) =\left| \breve{\Gamma}%
\left( f_{ref}\right) \right) \right\}  \nonumber \\
{\frak Co}_{S_{D}} &=&GL\left( S_{D}\right) /{\frak G}_{\Gamma \left(
f_{ref}\right) }
\end{eqnarray}%
One can see that the parameterization of $\ker \Xi $ given by ${\frak C}%
_{\ker \Xi }$, which does not depend on $f$ , (but it does implicitly if a
reference $\breve{Y}_{ref}$ is used as in Eq.$\left( \ref{ref}\right) $
rather than a reference $\breve{Y}_{ref\,K_{D}}$ that is totally contained
in $\ker \Xi $), and is provided by Eq.$\left( \ref{parker}\right) $.

\begin{itemize}
\item If $\phi \ $defines a vector bundle map then $\int P_{S_{D}}\left( 
{\sf z}_{1},{\sf z}_{1}^{\prime }\right) \phi \left( {\sf z}_{1}^{\prime
}\right) P_{K_{D}}\left( {\sf z}_{1}^{\prime },{\sf z}_{2}\right) d{\sf z}%
_{1}^{\prime }=0\quad \forall {\sf z}_{1},{\sf z}_{2}$

If $\phi \bot _{s}Y\quad \forall Y\in {\cal B}_{1}\left( {\cal H}^{N}\right) 
$ then $\Phi _{S_{D}}\left( {\sf z}_{1},{\sf z}_{2}\right) =\int
P_{S_{D}}\left( {\sf z}_{1},{\sf z}_{1}^{\prime }\right) \phi \left( {\sf z}%
_{1}^{\prime }\right) P_{S_{D}}\left( {\sf z}_{1}^{\prime },{\sf z}%
_{2}\right) d{\sf z}_{1}^{\prime }=0\quad \forall {\sf z}_{1},{\sf z}_{2}$

If $\phi \bot _{s}Y\quad \forall Y\in K_{D}$ then $\int P_{S_{D}}\left( {\sf %
z}_{1},{\sf z}_{1}^{\prime }\right) \phi \left( {\sf z}_{1}^{\prime }\right)
P_{K_{D}}\left( {\sf z}_{1}^{\prime },{\sf z}_{2}\right) d{\sf z}%
_{1}^{\prime }=0\quad \forall {\sf z}_{1},{\sf z}_{2}\Longleftrightarrow $ $%
\phi \ $defines a vector bundle map

If $\phi \not{\bot}_{s}Y$ then $\Phi _{S_{D}}\left( {\sf z}_{1},{\sf z}%
_{2}\right) =\int P_{S_{D}}\left( {\sf z}_{1},{\sf z}_{1}^{\prime }\right)
\phi \left( {\sf z}_{1}^{\prime }\right) P_{S_{D}}\left( {\sf z}_{1}^{\prime
},{\sf z}_{2}\right) d{\sf z}_{1}^{\prime }\neq 0\quad \forall {\sf z}_{1},%
{\sf z}_{2}$

If $\phi $ defines an element of ${\frak G}_{fib}$ then $\Phi _{S_{D}}\left( 
{\sf z}_{1},{\sf z}_{2}\right) =\int P_{S_{D}}\left( {\sf z}_{1},{\sf z}%
_{1}^{\prime }\right) \phi \left( {\sf z}_{1}^{\prime }\right)
P_{S_{D}}\left( {\sf z}_{1}^{\prime },{\sf z}_{2}\right) d{\sf z}%
_{1}^{\prime }=0\quad \forall {\sf z}_{1},{\sf z}_{2}$ and $\int
P_{K_{D}}\left( {\sf z}_{1},{\sf z}_{1}^{\prime }\right) \phi \left( {\sf z}%
_{1}^{\prime }\right) P_{S_{D}}\left( {\sf z}_{1}^{\prime },{\sf z}%
_{2}\right) d{\sf z}_{1}^{\prime }=0\quad \forall {\sf z}_{1},{\sf z}_{2}$

If $\phi $ defines an element of ${\frak G}_{den}$ then $\int
P_{K_{D}}\left( {\sf z}_{1},{\sf z}_{1}^{\prime }\right) \phi \left( {\sf z}%
_{1}^{\prime }\right) P_{S_{D}}\left( {\sf z}_{1}^{\prime },{\sf z}%
_{2}\right) d{\sf z}_{1}^{\prime }=0\quad \forall {\sf z}_{1},{\sf z}_{2}$
and $\Phi _{K_{D}}\left( {\sf z}_{1},{\sf z}_{2}\right) =\int
P_{K_{D}}\left( {\sf z}_{1},{\sf z}_{1}^{\prime }\right) \phi \left( {\sf z}%
_{1}^{\prime }\right) P_{K_{D}}\left( {\sf z}_{1}^{\prime },{\sf z}%
_{2}\right) d{\sf z}_{1}^{\prime }=0\quad \forall {\sf z}_{1},{\sf z}_{2}$

\item The subgroup ${\frak G}_{\ker }\subset {\frak G}_{\Xi }$ introduced in
the main section of paper I is formed by super multiplication operators
generated by congruence transformations composed of operators $\in {\cal B}%
_{1}\left( {\cal H}^{N}\right) $ that are diagonal in the coordinate
representation. The operators $\hat{T}\left( \phi \right) $ belonging to the
representation of ${\frak G}_{\ker }$ describe in the preceding have the
properties 
\begin{eqnarray}
\hat{T}\left( \phi ,\omega \right) &:&K_{OD}\rightarrow K_{OD}  \nonumber \\
\hat{T}\left( \phi ,\omega \right) &:&K_{D}\rightarrow K_{D}  \nonumber \\
\hat{T}\left( \phi ,\omega \right) &:&S_{D}\rightarrow K_{D}  \nonumber \\
\hat{T}\left( \phi ,\omega \right) &:&S_{D}\rightarrow S_{D}
\end{eqnarray}%
The subgroup that acts homogeneously on pure states is ${\frak G}_{pure}=$ $%
{\frak G}_{fib}\times {\frak G}_{den}\subset {\frak G}_{\ker }$. The
subgroup ${\frak G}_{\ker }$ can be factored as ${\frak G}_{\ker }={\frak Co}%
_{kfd}\times {\frak G}_{pure}$ where the coset ${\frak Co}_{kfd}{\frak =G}%
_{\ker }/\left( {\frak G}_{fib}\times {\frak G}_{den}\right) $ and contains
maps from $S_{D}\rightarrow K_{D}$. The group ${\frak G}_{\ker }$ can also
be factorized into products of positive and unitary operators as 
\begin{equation}
{\frak G}_{\ker }={\cal G}_{\ker }\times {\cal G}_{\Xi }
\end{equation}

\item The subgroup ${\frak G}_{den}$ can be represented by operators of the
form 
\begin{equation}
{\frak G}_{den}=\left( 
\begin{array}{ccc}
\hat{t}_{K_{OD}}\left( \hat{T}_{S_{D}}\right) & 0 & 0 \\ 
0 & \hat{t}_{K_{D}}\left( \hat{T}_{S_{D}}\right) & 0 \\ 
0 & 0 & \hat{T}_{S_{D}}%
\end{array}%
\right)
\end{equation}%
where $\hat{t}_{K_{OD}}$ and $\hat{t}_{K_{D}}$ are representations of $%
{\frak G}_{den}$. The representation $\hat{t}_{K_{OD}}$ is given by%
\begin{equation}
\hat{t}_{K_{OD}}\left( \hat{T}_{S_{D}}\right) \left( \breve{Y}%
_{K_{OD}}\right) =e^{{\cal \breve{V}+}i\breve{\omega}}\breve{Y}_{K_{OD}}e^{%
{\cal \breve{V}-}i\breve{\omega}}
\end{equation}%
for ${\cal V=V}\left( \phi _{S_{D}},\phi _{K_{D}0}\right) $ where $\phi
_{K_{D}0}$ and $\omega $ are fixed. Note that $\hat{T}_{S_{D}}$ is not a
super multiplication operator unless $\hat{t}_{K_{D}}$ is the identity
representation i.e. $\phi _{K_{D}0}=0$. The subgroup ${\frak G}_{fib}$ can
be represented by 
\begin{equation}
{\frak G}_{fib}=\left( 
\begin{array}{ccc}
\hat{t}_{K_{OD}}^{\prime }\left( \hat{T}_{K_{D}}\right) & 0 & 0 \\ 
0 & \hat{T}_{K_{D}} & 0 \\ 
0 & 0 & \hat{P}_{K_{S}}%
\end{array}%
\right)
\end{equation}%
where $\hat{t}_{K_{OD}}^{\prime }$ and $\hat{t}_{S_{D}}^{\prime }$ are
representations of ${\frak G}_{fib}$. The representation $\hat{t}%
_{K_{OD}}^{\prime }$ is given by 
\begin{equation}
\hat{t}_{K_{OD}}^{\prime }\left( \hat{T}_{K_{D}}\right) \left(
Y_{K_{OD}}\right) =e^{{\cal \breve{V}+}i\breve{\omega}}Y_{K_{OD}}e^{{\cal 
\breve{V}-}i\breve{\omega}}
\end{equation}%
for ${\cal V=V}\left( \phi _{K_{D}}\right) $ where $\omega $ is fixed.

\item Restricting attention to the subgroup ${\frak G}_{fib}$ of super
multiplication operators that transform equivalence classes to themselves
i.e. 
\begin{equation}
\hat{T}_{fib}:\left( \left| \breve{k}\right) ,\left| \Gamma \left( f\right)
\right) \right) \rightarrow \left( \left| \breve{k}^{\prime }\right) ,\left|
\Gamma \left( f\right) \right) \right) ;\quad \hat{T}_{fib}\in {\frak G}%
_{fib}
\end{equation}%
contained in $GL\left( \ker \Xi \right) $, leads to the non redundant
parameterization of pure states belonging to an equivalence class as
described in preceding. The parameterization for this action of ${\frak G}%
_{fib}$ on a reference element is given in Eq.$\left( \ref{refkd}\right) $
and the discussion preceding this equation, where the form of function, $%
\phi _{refK_{D}}$, generating the super operator transformation $\hat{T}%
_{fib}\left( \phi _{refK_{D}}\right) $ is described. The unitary subgroup $%
{\cal G}_{\Xi }$ is contained in ${\frak G}_{fib}$ and we have used the
factorization ${\frak G}_{fib}={\cal G}_{fib}\times {\cal G}_{\Xi }$
extensively in section V of paper I. The function $\phi _{refK_{D}}$ is
strongly orthogonal to $Y$, the integral kernel of any pure state, as by
construction $\hat{P}_{S_{D}}\hat{\Phi}_{refK_{D}}\hat{P}_{S_{D}}=0$ (see
Eqs.$\left( \ref{diagker}\right) $,$\left( \ref{kerproj}\right) $and $\left( %
\ref{PhiKD}\right) )$ thus has no affect on $\hat{P}_{S_{D}}Y$ for any $Y$,
while $\hat{\Phi}_{refK_{D}}:K_{D}\rightarrow K_{D}$ and $\hat{\Phi}%
_{refK_{OD}}:K_{OD}\rightarrow K_{OD}$, which means that their integral
kernels are automatically strongly orthogonal to $P_{K_{OD}}Y$ and $%
P_{K_{D}}Y$. {\it Note }that it is possible for a more general $\phi $ (that
has non zero components, $P_{S_{D}}\phi )$, to be strongly orthogonal to a
specific $P_{S_{D}}Y_{ref}$ but not to {\it all }$P_{S_{D}}Y$.

\item Restricting attention to the subgroup ${\frak G}_{den}$ of super
multiplication operators contained in $GL\left( S_{D}\right) $, represented
in the form of Eq.$\left( \ref{tden}\right) $,\ leads to the non redundant
parameterization of a path pure states corresponding to different densities
in a 1-1 fashion, (see Eqs.$\left( \ref{tden}\right) $,$\left( \ref{tden2}%
\right) $and $\left( \ref{tden3}\right) )$.
\end{itemize}

\subsection{Example\protect\bigskip}

As an example we consider the pure Independent Particle State $\breve{Y}%
_{ref}\left( \gamma _{0}\right) =\left| \Psi _{0\,IPS}\right\rangle
\left\langle \Psi _{0\,IPS}\right| $, where $\Xi \left( \breve{Y}%
_{ref}\left( \gamma _{0}\right) \right) =\gamma _{0}$, which has the
decomposition 
\begin{equation}
\breve{Y}_{ref}=\breve{Y}_{ref\,K_{OD}}+\breve{Y}_{ref\,K_{D}}+\breve{\Gamma}%
\left( \gamma _{0}\right)
\end{equation}%
{\bf All} pure states with the same density as this reference state are
generated by 
\begin{equation}
\left| \breve{Y}\left( \gamma _{0}\right) \right) =\hat{T}_{fib}\left( \phi
_{K_{D}}{}_{ref}\right) \left| \breve{Y}_{ref}\left( \gamma _{0}\right)
\right) ;\quad \hat{T}_{fib}\left( \phi _{K_{D}}{}_{ref}\right) \in {\frak G}%
_{fib}
\end{equation}%
A {\it full }path of path pure states corresponding to {\bf all} different
densities in a 1-1 fashion is generated by 
\begin{equation}
\left| \breve{Y}_{ref}\left( \gamma \right) \right) =\hat{T}_{den}\left(
\phi _{\Gamma \left( \gamma _{0}\right) }\right) \left| \breve{Y}%
_{ref}\left( \gamma _{0}\right) \right) ;\,\gamma \neq \gamma ^{\prime
},\quad \hat{T}_{den}\left( \phi _{\Gamma \left( \gamma _{0}\right) }\right)
\in {\frak G}_{den}
\end{equation}%
{\bf All }pure states with the density $\gamma $ are generated from $\left| 
\breve{Y}_{ref}\left( \gamma \right) \right) $ by%
\begin{equation}
\left| \breve{Y}\left( \gamma \right) \right) =\hat{T}_{fib}\left( \phi
_{K_{D}}{}_{ref\left( \gamma \right) }\right) \left| \breve{Y}_{ref}\left(
\gamma \right) \right) ;\quad \hat{T}_{fib}\left( \phi _{K_{D}}{}_{ref\left(
\gamma \right) }\right) \in {\frak G}_{fib}
\end{equation}%
A full path can be produced by finding intra fiber transformations $\hat{T}%
_{fib\,IPS}\left( \phi _{S_{D}}\left( \gamma \right) \right) $, fixed for
each fiber, that produce an IPS\ from $\left| \breve{Y}\left( \gamma \right)
\right) $ 
\begin{equation}
\left| \breve{Y}_{IPS}\left( \gamma \right) \right) =\hat{T}%
_{fib\,IPS}\left( \phi _{S_{D}}\left( \gamma \right) \right) \left| \breve{Y}%
\left( \gamma \right) \right) ;\quad \hat{T}_{fib\,IPS}\left( \phi
_{S_{D}}\left( \gamma \right) \right) \in {\frak G}_{fib}
\end{equation}%
The transformations $\hat{T}_{fib\,IPS}\left( \left( \phi _{S_{D}}\left(
\gamma \right) \right) \right) $ are not unique as in general there many
IPS\ in any given fiber/equivalence class. Thus one can produce many full
paths of IPS.

\begin{references}
\bibitem{ideal} A two sided ideal of ${\cal B}\left( {\cal H}^{N}\right) $
is a set ${\cal Y\subseteq B}\left( {\cal H}^{N}\right) $ such that $YX\in 
{\cal B}\left( {\cal H}^{N}\right) $ and $XY\in {\cal B}\left( {\cal H}%
^{N}\right) \forall X\in {\cal B}\left( {\cal H}^{N}\right) $ and $\forall
Y\in {\cal Y}$.

\bibitem{homogen} A groups of transformations acts homogeneously on a space
if fro every $Y_{1}$ and $Y_{2}$ that belong to the space there exists a
group element $g$ such that $Y_{1}=gY_{2}$.
\end{references}
